% =====================================================================
%  HISTORIAL TÉCNICO COMPLETO DEL PROYECTO
%  Ayuda-memoria de todo lo investigado, desarrollado, probado,
%  descartado y seleccionado — de febrero a agosto de 2026.
%
%  Este NO es el documento de entrega. La entrega formal (Primera
%  Presentación, Art. 19, máx. 15 páginas) vive en ../latex/.
%  Este archivo existe para recordar el trabajo hecho, con datos y
%  figuras reales, sin límite de extensión.
%
%  Autor: Elías David Álvarez Martínez
%  Compilar: pdflatex main -> pdflatex main
% =====================================================================
\documentclass[11pt,a4paper]{article}

\usepackage[utf8]{inputenc}
\usepackage[T1]{fontenc}
\usepackage[spanish,es-tabla,es-noquoting]{babel}
\usepackage{lmodern}
\usepackage[left=2.2cm,right=2.2cm,top=2.4cm,bottom=2.2cm]{geometry}
\usepackage{amsmath,amssymb}
\usepackage{graphicx}
\usepackage{array,tabularx,booktabs,multirow,longtable}
\usepackage[table]{xcolor}
\usepackage{enumitem}
\usepackage{titlesec}
\usepackage{microtype}
\usepackage{siunitx}
\usepackage{caption}
\usepackage[hidelinks]{hyperref}

\sisetup{output-decimal-marker={,},per-mode=symbol}
\graphicspath{%
  {figuras/}%
  {../../Urucom_2026_compact/imagenes/}%
  {../../calibration_fsm/}%
  {../../investigacion/Assets/}%
}
\captionsetup{font=small,labelfont=bf}

\setlength{\parindent}{0pt}
\setlength{\parskip}{0.45em}

% Rutas y nombres de archivo largos en tablas: permitir corte y dar aire al
% algoritmo de linea antes de que desborde el margen.
\setlength{\emergencystretch}{3em}
\sloppy
\makeatletter
\renewcommand{\ttdefault}{lmtt}
\makeatother

% --- Encabezado manual (fancyhdr NO está instalado en este MiKTeX) ---
\makeatletter
\renewcommand{\@oddhead}{%
  \parbox{\textwidth}{%
    \footnotesize\itshape Historial t\'ecnico --- ayuda memoria del proyecto
    \hfill \upshape\thepage\\[-0.6em]\rule{\textwidth}{0.4pt}}}
\renewcommand{\@evenhead}{\@oddhead}
\renewcommand{\@oddfoot}{}
\renewcommand{\@evenfoot}{}
\makeatother

% --- Cajas semánticas -------------------------------------------------
\definecolor{cajaDesc}{HTML}{FBEAEA}   % descartado
\definecolor{cajaSel}{HTML}{E8F2EA}    % seleccionado / conclusión
\definecolor{cajaDato}{HTML}{EAEFF7}   % dato duro / medición
\definecolor{cajaDuda}{HTML}{F2F2F2}   % pendiente / a confirmar

\newcommand{\cajita}[3]{%
  \par\vspace{0.35em}%
  \noindent\fcolorbox{black!25}{#1}{%
    \parbox{\dimexpr\linewidth-2\fboxsep-2\fboxrule}{\footnotesize\textbf{#2} #3}}%
  \par\vspace{0.35em}}

\newcommand{\descartado}[1]{\cajita{cajaDesc}{[Descartado]}{#1}}
\newcommand{\seleccionado}[1]{\cajita{cajaSel}{[Se qued\'o]}{#1}}
\newcommand{\dato}[1]{\cajita{cajaDato}{[Dato medido]}{#1}}
\newcommand{\pendiente}[1]{\cajita{cajaDuda}{[Pendiente / a confirmar]}{#1}}

% Etiqueta de rango de fechas al inicio de cada etapa
\newcommand{\periodo}[1]{{\small\bfseries\itshape #1}\par\vspace{0.2em}}

\begin{document}

% ---------------------------------------------------------------- PORTADA
\thispagestyle{empty}
\begin{center}
    \vspace*{1.5cm}
    {\Large\bfseries Historial t\'ecnico del proyecto}\par
    \vspace{0.5cm}
    {\large Dise\~no e implementaci\'on de un sistema electr\'onico multicanal\\
    del tipo IoT para caracterizaci\'on geot\'ecnica del suelo\\
    mediante ondas s\'ismicas}\par
    \vspace{1.0cm}
    {\itshape Ayuda-memoria de todo lo investigado, desarrollado, probado,\\
    descartado y seleccionado}\par
    \vspace{0.3cm}
    {25 de febrero --- 12 de agosto de 2026}\par
    \vspace{1.2cm}
    \begin{tabularx}{0.8\textwidth}{@{}lX@{}}
    \textbf{Alumno} & El\'ias David \'Alvarez Mart\'inez (Y24137) \\
    \textbf{Carrera} & Ingenier\'ia Electr\'onica \\
    \textbf{Tutor} & Enrique A. Vargas Cabral, Ph.D. \\
    \textbf{Instituci\'on} & UCA --- Facultad de Ciencias y Tecnolog\'ias \\
    \end{tabularx}
\end{center}

\vspace{1.0cm}
\hrule
\vspace{0.5cm}

\noindent\textbf{Qu\'e es este documento.} No es la entrega de la Primera
Presentaci\'on ---esa es un documento de Estado del Arte de m\'aximo 15
p\'aginas y vive en \texttt{docs/Primera Presentaci\'on/latex/}. Esto es el
registro completo del trabajo: qu\'e se ley\'o, qu\'e se dise\~n\'o, qu\'e se
midi\'o, qu\'e fall\'o y por qu\'e se descart\'o. Est\'a reconstruido a partir de
la bit\'acora diaria (66 entradas, \textasciitilde{}1{,}3~MB), del historial de
\emph{commits} de los diez repositorios del proyecto (1262 \emph{commits}) y de
los datos crudos y procesados de las campa\~nas de campo.

\vspace{0.3cm}
\noindent\textbf{C\'omo leerlo.} Est\'a ordenado cronol\'ogicamente en etapas.
Dentro de cada etapa se distingue con color lo que se
\colorbox{cajaSel}{\strut se qued\'o}, lo que se
\colorbox{cajaDesc}{\strut se descart\'o}, los
\colorbox{cajaDato}{\strut datos medidos} y lo que queda
\colorbox{cajaDuda}{\strut pendiente}.

\newpage
\tableofcontents
\newpage

\section{Panorama general}
\label{sec:panorama}

\subsection{De qu\'e se trata el proyecto}

El objetivo es construir un sismógrafo multicanal de bajo costo para
caracterizaci\'on geot\'ecnica del suelo mediante el m\'etodo MASW
(\emph{Multichannel Analysis of Surface Waves}): se golpea el suelo con una
fuente de impacto, se registra la llegada de las ondas Rayleigh en un arreglo
lineal de ge\'ofonos, se construye la imagen de dispersi\'on
velocidad-de-fase/frecuencia y se invierte esa curva para obtener el perfil de
velocidad de onda de corte $V_S(z)$.

Lo que distingue a este sistema de un sism\'ografo comercial es el conjunto de
restricciones autoimpuestas:

\begin{itemize}[itemsep=0.2em]
\item \textbf{El transductor es el que hay}, no el que conviene: ge\'ofonos
SM-24 de \SI{10}{\hertz} disponibles en la Facultad. Para llegar a las
profundidades de inter\'es hace falta informaci\'on por debajo de su frecuencia
natural, de modo que \emph{extender la banda hacia abajo por v\'ia electr\'onica
deja de ser un refinamiento y pasa a ser una condici\'on de dise\~no}. La
extensi\'on efectiva todav\'ia debe medirse de extremo a extremo.
\item \textbf{Nodos inal\'ambricos e independientes}, no un cable de decenas de
metros con un digitalizador central. Eso convierte la sincronizaci\'on entre
canales en un problema de dise\~no de primer orden.
\item \textbf{Orquestaci\'on y revisi\'on en herramientas propias}, desde la
adquisici\'on hasta una fachada de inversi\'on con motores propios y de terceros.
La capacidad multimodal existe, pero el estado persistido usa s\'olo una curva
etiquetada como modo~0, sin identificaci\'on independiente como fundamental.
\end{itemize}

\subsection{Resumen autoritativo al 12 de agosto de 2026}

Esta tabla es el punto de entrada para resumir el proyecto. Ante una diferencia
con un relato cronol\'ogico posterior, prevalecen los archivos de resultados, el
firmware vigente y los metadatos indicados aqu\'i.

\begin{center}\footnotesize
\begin{tabularx}{\textwidth}{@{}l X@{}}
\toprule
\textbf{Eje} & \textbf{Estado que puede afirmarse} \\
\midrule
Transductor & SM-24 disponible en la Facultad: $f_n=\SI{10}{\hertz}$,
sensibilidad nominal \SI{28.8}{\volt\second\per\meter}, $\zeta_0=0{,}25$,
bobina \SI{375}{\ohm} y masa m\'ovil \SI{11}{\gram}; la hoja de datos no
especifica inductancia. \\
\addlinespace[2pt]
Compensaci\'on & En la tanda calibrada, $\zeta_0=0{,}25$ es el objetivo del
numerador y $\zeta_1=937{,}396$ el denominador BP. $83{,}661$ pertenece a la
tanda descartada. El polo inferior ideal es \SI{5.44}{\milli\hertz}, pero la
muesca medida qued\'o \SI{26.33}{\decibel} menos profunda que la de dise\~no. \\
\addlinespace[2pt]
Modelo del operacional & Los scripts combinaron $A_0=\SI{90}{\decibel}$ con
\SI{8}{\mega\hertz} como polo dominante. Si este valor es GBW, el polo monopolo
es aproximadamente \SI{253}{\hertz}. La relaci\'on ya se corrigi\'o en los scripts;
la s\'intesis y el contraste deben repetirse y el ajuste previo no se usa como validaci\'on. \\
\addlinespace[2pt]
DA-AFE & Sobre dos placas y ocho configuraciones, la salida LP pas\'o de
\SIrange{2.9}{45.8}{\percent} a \SIrange{0.4}{3.9}{\percent} del fondo de
escala, con 10 reinicios apareados por configuraci\'on y un \'unico PI. \\
\addlinespace[2pt]
ADC y adquisici\'on & Firmware actual: \SI{2604}{\hertz} nativos, 18 bits,
$F_s=2604/N$ con $N=1,\ldots,100$ y cuatro rangos. Las tasas 2929 y
\SI{1020}{\hertz} son hist\'oricas de campo. La ruta cr\'itica usa DFB, DMA y
\texttt{superMaquina} en UDB. \\
\addlinespace[2pt]
Sincronizaci\'on & PRESTART/HOT\_WAIT asegura disponibilidad; el
\texttt{CMD\_START} llega por ESP-NOW y cada esclavo genera el flanco local al
PSoC. La observaci\'on $<\SI{400}{\micro\second}$ a unos \SI{5}{\meter} es
preliminar y anterior a la medici\'on completa de la topolog\'ia vigente. \\
\addlinespace[2pt]
Campo & Apertura sint\'etica: martillo fijo en $x=0$ y un ge\'ofono m\'ovil. El
grupo principal tiene 21 posiciones entre \SI{10}{} y \SI{50}{\meter},
$\Delta x=\SI{2}{\meter}$ y $L=\SI{40}{\meter}$. El orden de profundidad con
soporte conservador es \SI{20}{\meter}, no \SI{50}{\meter}. \\
\addlinespace[2pt]
MASW & El exportado de julio tiene 51 puntos y desajuste declarado de
\SI{2.554}{\percent}, pero llega a $\lambda=\SI{130.6}{\meter}$ y queda fuera de
la evidencia. El estado posterior persistido usa $\lambda=\SIrange{4.001}{26.732}{\meter}$,
\texttt{maswavespy(port)}, modo~0, siete capas y desajuste \SI{7.1063}{\percent}. No hay
validaci\'on SPT ni resultado multimodal comparado. \\
\addlinespace[2pt]
Pendientes de cierre & Modelo no ideal del operacional, banda compensada medida, \emph{jitter} extremo a
extremo, apareamiento entre canales, mayor apertura y energ\'ia de fuente,
ruido referido a entrada y contraste geot\'ecnico externo. \\
\bottomrule
\end{tabularx}
\end{center}

\begin{figure}[!ht]
\centering
\includegraphics[width=\textwidth]{arquitectura_vigente.png}
\caption{Arquitectura vigente y frontera entre el camino determinista de
muestra y las funciones de calibraci\'on, comunicaci\'on, almacenamiento y
procesamiento.}
\label{fig:arquitectura-vigente-hist}
\end{figure}

\subsection{Escala del trabajo}

\begin{figure}[!ht]
\centering
\includegraphics[width=\textwidth]{timeline_commits.png}
\caption{Actividad de desarrollo semanal en los diez repositorios del proyecto.
La intensidad codifica cantidad de \emph{commits} por semana en escala
comprimida (ra\'iz cuadrada), porque de otro modo la sesi\'on de 233
\emph{commits} de abril aplastar\'ia todo el resto. La fila
\emph{Superproyecto} concentra la historia anterior a la modularizaci\'on de
julio de 2026: hasta ese momento todo viv\'ia en un \'unico repositorio, y por
eso las filas de los subm\'odulos arrancan m\'as tarde aunque el trabajo sea
anterior.}
\label{fig:timeline}
\end{figure}

\begin{center}\small
\begin{tabularx}{\textwidth}{@{}lX r@{}}
\toprule
\textbf{Dimensi\'on} & \textbf{Detalle} & \textbf{Cantidad} \\
\midrule
Per\'iodo activo & 25 feb 2026 --- 12 ago 2026 & 25 semanas \\
\emph{Commits} & suma reconstruida de los diez repositorios & \textasciitilde{}1261 \\
D\'ias con actividad & registrados en la bit\'acora & 66 \\
Bit\'acora & entradas diarias con an\'alisis de \emph{diff} & \textasciitilde{}1{,}3~MB \\
Bibliograf\'ia & base curada con DOI y relevancia & 65 papers \\
Notas de concepto & vault Obsidian, conceptos at\'omicos & \textasciitilde{}153 \\
Datos crudos de campo & tres sitios, LFS con dedup SHA-256 & \textasciitilde{}29~GiB \\
Datos procesados & \emph{waterfalls}, dispersi\'on, perfiles & \textasciitilde{}1{,}1~GiB \\
\bottomrule
\end{tabularx}
\end{center}

\subsection{Los diez repositorios}

En julio de 2026 el monolito se parti\'o en un superproyecto con subm\'odulos
organizados \emph{por prop\'osito}, no por lenguaje. La organizaci\'on final es:

\begin{center}\small
\begin{tabularx}{\textwidth}{@{}l l X@{}}
\toprule
\textbf{Ruta} & \textbf{Repositorio} & \textbf{Contenido} \\
\midrule
\texttt{src/firmware/psoc} & Tesis-firmware-psoc & Adquisici\'on: cadena anal\'ogica configurable, ADC $\Delta\Sigma$, DFB, DMA, m\'aquina de estados en UDB, FatFs sobre SD \\
\texttt{src/firmware/esp32} & Tesis-firmware-esp32 & Red ESP-NOW maestro/esclavo, sincronizaci\'on de disparo, volcado de capturas, servidor web embebido \\
\texttt{src/interfaces/python} & Tesis-interfaces-python & Servidor FastAPI, interfaz web de revisi\'on de campo, cola de an\'alisis \\
\texttt{src/interfaces/matlab} & Tesis-interfaces-matlab & Osciloscopio virtual y bancos de prueba contra el hardware \\
\texttt{src/calculos\_modelados/matlab} & Tesis-calculos-matlab & Modelo del SM-24, compensadores, identificaci\'on de la cadena anal\'ogica, martillo de leva en Simscape \\
\texttt{src/calculos\_modelados/python} & Tesis-calculos-python & C\'alculo de compensadores, MASW, martillo de leva 2D \\
\texttt{docs} & Tesis-documentacion & Vault de investigaci\'on, bit\'acora, papers, art\'iculo URUCON, esta presentaci\'on \\
\texttt{data} & Tesis-datos & Registros crudos y procesados de campo y de osciloscopio \\
\texttt{PCBs} & Tesis-PCBs & Dise\~nos de placa \\
\midrule
\multicolumn{3}{@{}l@{}}{\emph{(el superproyecto \texttt{Tesis} sostiene los \emph{gitlinks} y la historia previa a la partici\'on)}} \\
\bottomrule
\end{tabularx}
\end{center}

\subsection{Las etapas}

\begin{center}\small
\begin{tabularx}{\textwidth}{@{}l l X@{}}
\toprule
\textbf{\S} & \textbf{Per\'iodo} & \textbf{Qu\'e pas\'o} \\
\midrule
\ref{sec:et1} & 25 feb -- 23 mar & Marco te\'orico: Foti completo, 65 papers, scripts de validaci\'on \\
\ref{sec:et2} & 10 -- 15 abr & Investigaci\'on asistida: 140 conceptos at\'omicos \\
\ref{sec:et3} & 17 abr -- 20 may & Electr\'onica anal\'ogica: modelo del SM-24, compensaci\'on, dos pivotes de topolog\'ia \\
\ref{sec:et4} & 24 may -- 3 jun & Red ESP-NOW: sincronizaci\'on de disparo entre nodos \\
\ref{sec:et5} & 4 -- 10 jun & Consolidaci\'on del nodo: firmware \'unico e interfaz web embebida \\
\ref{sec:et6} & 14 jun -- 2 jul & Auto-calibraci\'on, DMA y \texttt{superMaquina} \\
\ref{sec:et7} & 3 -- 9 jul & Campa\~nas de campo y \emph{pipeline} MASW en Python \\
\ref{sec:et8} & 11 -- 12 jul & Auditor\'ia pre-campo, SD en el PSoC, tres bugs cr\'iticos \\
\ref{sec:et9} & 15 -- 23 jul & Art\'iculo URUCON y caracterizaci\'on medida de la cadena \\
\ref{sec:et10} & 21 jul -- 1 ago & Ruptura arquitect\'onica y portaci\'on de la interfaz a web \\
\ref{sec:et11} & 2 -- 12 ago & Placa nueva: rama \texttt{cambios-hardware} \\
\bottomrule
\end{tabularx}
\end{center}

\section{Etapa 1 --- Marco te\'orico}
\label{sec:et1}
\periodo{25 de febrero --- 23 de marzo de 2026 \quad$\cdot$\quad 22 \emph{commits} en 9 d\'ias de trabajo}

Esta etapa no produjo hardware ni c\'odigo de producci\'on. Produjo el
\emph{criterio} con el que se tomaron todas las decisiones posteriores, y por
eso vale la pena tenerla registrada: buena parte de lo que despu\'es se
descart\'o se descart\'o por razones que se aprendieron ac\'a.

\subsection{Punto cero}

El repositorio se cre\'o el 25 de febrero de 2026 a las 16:24. Un solo
\emph{commit}, sin contenido t\'ecnico. La primera sesi\'on real fue el 3 de
marzo: cuatro horas continuas en las que se mont\'o el vault de Obsidian y se
escribi\'o el primer resumen acad\'emico del Cap\'itulo~1 de Foti. De ah\'i sali\'o
la metodolog\'ia que gobern\'o toda la investigaci\'on posterior:

\begin{itemize}[itemsep=0.2em]
\item un archivo por cap\'itulo con el resumen extenso;
\item \textbf{conceptos at\'omicos} enlazados: una nota por idea, con
definici\'on, intuici\'on f\'isica y criterio cuantitativo cuando existe;
\item enlaces bidireccionales entre conceptos, de modo que el grafo del vault
muestre qu\'e ideas sostienen a cu\'ales.
\end{itemize}

Del Cap\'itulo~1 salieron los 9 primeros conceptos at\'omicos. Al cerrar la
etapa hab\'ia m\'as de 90, y tras la etapa siguiente superaron los 150.

\subsection{Fuente primaria: Foti \emph{et al.}, los ocho cap\'itulos}

La referencia estructural fue \emph{Surface Wave Methods for Near-Surface Site
Characterization} de Sebastiano Foti y colaboradores. Se resumieron los ocho
cap\'itulos completos entre el 3 y el 19 de marzo. Lo que qued\'o de cada uno,
en t\'erminos de lo que despu\'es se us\'o:

\begin{center}\small
\begin{tabularx}{\textwidth}{@{}l l X@{}}
\toprule
\textbf{Cap.} & \textbf{Tema} & \textbf{Qu\'e se us\'o despu\'es} \\
\midrule
1 & Panorama del m\'etodo &
Problema directo vs.\ inverso; el flujo adquisici\'on\,$\to$\,dispersi\'on\,$\to$\,inversi\'on que estructura todo el sistema; comparaci\'on $V_{S30}$ invasivo vs.\ ondas superficiales \\
\addlinespace[2pt]
2 & Propagaci\'on en medios estratificados &
Relaci\'on de dispersi\'on, tensores de esfuerzo y deformaci\'on, Ley de Hooke generalizada, derivaci\'on por potenciales el\'asticos, problema de Lamb (1904) \\
\addlinespace[2pt]
3 & Medici\'on &
\textbf{El cap\'itulo de dise\~no experimental.} Traza s\'ismica y \emph{gather} multicanal, transformada 2D al dominio $(k,f)$, Nyquist en tiempo \emph{y en espacio} ($\lambda_{\min}=2\Delta x$), reglas de \emph{offset} m\'inimo y longitud de arreglo \\
\addlinespace[2pt]
4--5 & Dispersi\'on y atenuaci\'on &
Velocidad de fase vs.\ de grupo; SASW y MOPA como antecedentes del multicanal \\
\addlinespace[2pt]
6 & Inversi\'on &
\emph{Ill-posedness} y no unicidad; b\'usqueda local vs.\ global. Es la justificaci\'on de por qu\'e despu\'es se eligi\'o un algoritmo evolutivo y no un gradiente \\
\addlinespace[2pt]
7 & Casos de estudio & Contraste de \'ordenes de magnitud esperables \\
\addlinespace[2pt]
8 & M\'etodos avanzados &
Ondas Love, ondas de Scholte, interferometr\'ia s\'ismica a partir de ruido ambiental \\
\bottomrule
\end{tabularx}
\end{center}

Dos resultados del Cap\'itulo~3 se convirtieron en restricciones duras del
dise\~no y conviene tenerlos a mano:

\begin{itemize}[itemsep=0.2em]
\item \textbf{\emph{Aliasing} espacial.} El espaciamiento entre ge\'ofonos
$\Delta x$ acota el n\'umero de onda recuperable: $k_{\max}=1/(2\Delta x)$, es
decir $\lambda_{\min}=2\Delta x$. Violar esto no degrada suavemente: hace que
los modos r\'apidos aparezcan como modos lentos en el espectro, o sea produce
un perfil $V_S$ \emph{plausible pero equivocado}.
\item \textbf{Efecto de campo cercano.} A distancias cortas de la fuente el
campo de ondas todav\'ia no est\'a dominado por la onda Rayleigh plana, y la
velocidad de fase estimada se sesga. Fija un \emph{offset} m\'inimo.
\end{itemize}

\pendiente{Estos dos criterios son exactamente los que no se pueden verificar
todav\'ia contra las campa\~nas realizadas, porque falta registrar la geometr\'ia
usada (espaciamiento, longitud del arreglo, \emph{offset} m\'inimo). Ver
\S\ref{sec:pendientes}.}

\subsection{Del papel al c\'odigo: los cinco scripts de validaci\'on te\'orica}

El 17 y 18 de marzo el proyecto cruz\'o de ``leer y resumir'' a ``programar''.
Se incorpor\'o como \emph{submodule} la librer\'ia MASW-Matlab de terceros y se
escribieron cinco scripts propios cuyo prop\'osito era comprobar que la teor\'ia
le\'ida se entend\'ia de verdad ---es decir, reproducirla num\'ericamente:

\begin{center}\small
\begin{tabularx}{\textwidth}{@{}l X@{}}
\toprule
\textbf{Script} & \textbf{Qu\'e verifica} \\
\midrule
\texttt{p01\_ondas\_de\_cuerpo.m} & Propagaci\'on P y S, relaciones con los par\'ametros de Lam\'e \\
\texttt{p02\_rayleigh\_love.m} & Existencia y forma de los modos superficiales \\
\texttt{p03\_sismograma\_imagen\_dispersion.m} & \textbf{El puente al sistema real}: de un \emph{gather} sint\'etico a la imagen de dispersi\'on \\
\texttt{p04\_inversion\_perfil\_vs.m} & Recuperaci\'on del perfil $V_S$ a partir de la curva \\
\texttt{p05\_viscoelasticidad.m} & Atenuaci\'on y su efecto sobre el espectro \\
\bottomrule
\end{tabularx}
\end{center}

\seleccionado{El flujo del \texttt{p03} ---\emph{gather} sint\'etico
$\to$ transformada $\to$ imagen de dispersi\'on--- es el mismo que despu\'es se
reimplement\'o en Python en \texttt{masw\_dispersion.py} (Etapa~7). La
versi\'on MATLAB sigui\'o siendo \'util como referencia contra la cual contrastar
la implementaci\'on propia.}

\subsection{La base bibliogr\'afica: de 22 a 65 papers}

Entre el 20 y el 23 de marzo se construy\'o la base de datos de literatura.
Result\'o en una tabla estructurada con DOI, categor\'ia, relevancia declarada
para la tesis y observaciones ---no una lista de t\'itulos--- que vive en
\texttt{docs/investigacion/sources/research/tables/} y de la cual sali\'o
despu\'es la bibliograf\'ia del documento de entrega.

\begin{center}\small
\begin{tabularx}{\textwidth}{@{}l X@{}}
\toprule
\textbf{Referencia} & \textbf{Por qu\'e qued\'o en la base} \\
\midrule
Park, Miller \& Xia (1999) & El paper fundacional del MASW: define el m\'etodo multicanal y el \emph{phase-shift} \\
Aki (1957) & Origen de los m\'etodos pasivos (SPAC); base te\'orica de la extensi\'on a ruido ambiental \\
Foti \emph{et al.} (2018), InterPACIFIC & Comparaci\'on ciega entre laboratorios: cu\'anta dispersi\'on entre operadores es normal. Es la vara con la que medirse \\
Moffat, Correia \& Past\'en (2016) & Validaci\'on en contexto latinoamericano \\
Alhuay-Le\'on \& Trejo-Nore\~na (2021) & Correlaci\'on $V_S$--SPT en Per\'u: el camino de validaci\'on m\'as accesible ac\'a \\
SESAME (2004) & Gu\'ia de referencia para HVSR \\
\bottomrule
\end{tabularx}
\end{center}

Un grupo de la base merece menci\'on aparte porque orient\'o directamente el
dise\~no: los papers de \textbf{hardware de bajo costo}. Son los que
establecieron que el problema no era conseguir un ADC bueno sino resolver la
banda baja del transductor y la sincronizaci\'on entre nodos.

\subsection{Balance de la etapa}

\seleccionado{\textbf{Qued\'o:} la estructura de conceptos at\'omicos como forma
de trabajo; la librer\'ia MASW-Matlab como referencia de contraste; los cinco
scripts de validaci\'on; la base de 65 papers con DOI y relevancia; y los dos
criterios de dise\~no de arreglo (\emph{aliasing} espacial y campo cercano).}

\descartado{\textbf{No qued\'o:} nada se descart\'o formalmente en esta etapa
---no hab\'ia todav\'ia qu\'e descartar. Lo que s\'i qued\'o desactualizado fue la
organizaci\'on de carpetas del vault, que se reescribi\'o dos veces (13 y 17 de
marzo) antes de estabilizarse en la taxonom\'ia numerada 00--08.}

\section{Etapa 2 --- Investigaci\'on asistida}
\label{sec:et2}
\periodo{10 --- 15 de abril de 2026 \quad$\cdot$\quad \textasciitilde{}119 \emph{commits} en 4 d\'ias}

Etapa corta en calendario y enorme en volumen. Se defini\'o un procedimiento de
trabajo ---en la jerga de la herramienta, una \emph{skill} llamada
\texttt{INVESTIGADOR}--- y se lo dej\'o correr de forma aut\'onoma sobre el vault
de investigaci\'on durante dos d\'ias.

\begin{figure}[ht]
\centering
\includegraphics[width=0.96\textwidth]{comparacion_metodos.png}
\caption{S\'intesis visual del espacio de m\'etodos que organiz\'o la investigaci\'on
asistida. Resume funci\'on y log\'istica; no establece una jerarqu\'ia universal.}
\end{figure}

\subsection{Qu\'e se le pidi\'o al agente}

El procedimiento fijaba cuatro cosas, y vale la pena registrarlas porque son
las que hicieron que el resultado fuera utilizable y no ruido:

\begin{enumerate}[itemsep=0.2em]
\item \textbf{Naturaleza anal\'itica}: no resumir por resumir, sino explicar el
mecanismo f\'isico.
\item \textbf{Gesti\'on de fuentes}: todo lo que se afirme debe rastrearse a la
base de 65 papers ya construida en la Etapa~1 ---sin fuente, no entra.
\item \textbf{Arquitectura en Obsidian}: un concepto por archivo, con bloque de
definici\'on, ejemplo concreto tomado de un paper real, y enlaces salientes.
\item \textbf{Criterio cient\'ifico}: distinguir lo establecido de lo
especulativo.
\end{enumerate}

\subsection{Qu\'e produjo}

\begin{center}\small
\begin{tabularx}{\textwidth}{@{}X r@{}}
\toprule
\textbf{Resultado} & \textbf{Cantidad} \\
\midrule
Conceptos at\'omicos enriquecidos (definici\'on + ejemplo + enlaces) & 65 \\
Conceptos at\'omicos nuevos & \textasciitilde{}37 \\
Total de conceptos en el vault al cierre & \textasciitilde{}140 \\
Ejemplos concretos por cap\'itulo (uno por cada uno de los 8) & 4 c/u \\
\emph{Wikilinks} rotos al cerrar la verificaci\'on & 0 \\
Archivos hu\'erfanos (sin ning\'un enlace entrante) & 0 \\
\bottomrule
\end{tabularx}
\end{center}

Los conceptos nuevos se agruparon en tres bloques que anticipan exactamente lo
que despu\'es hubo que implementar:

\begin{itemize}[itemsep=0.2em]
\item \textbf{Procesamiento de se\~nales}: transformadas, ventaneo,
\emph{phase-shift} ---base del \texttt{masw\_dispersion.py} de la Etapa~7.
\item \textbf{Inversi\'on avanzada}: no unicidad, regularizaci\'on, b\'usqueda
global ---justificaci\'on del uso de un algoritmo evolutivo en la Etapa~7.
\item \textbf{Se\~nales discretas}: muestreo, \emph{aliasing}, cuantizaci\'on
---base del dise\~no del ADC y la decimaci\'on en la Etapa~3.
\end{itemize}

\subsection{Consolidaci\'on y cierre}

El 15 de abril se revis\'o manualmente lo producido, se reorganizaron las
im\'agenes en \texttt{Assets/} y se cerr\'o la rama. Este paso importa: el
resultado del agente no se acept\'o autom\'aticamente, se revis\'o antes de
integrarlo.

\seleccionado{\textbf{Qued\'o:} el vault de \textasciitilde{}140 conceptos
enlazados, que funcion\'o despu\'es como referencia r\'apida durante todo el
dise\~no ---sobre todo los conceptos de procesamiento y de inversi\'on. Tambi\'en
qued\'o el m\'etodo: definir el criterio primero, dejar correr, revisar despu\'es.}

\pendiente{El vault no se volvi\'o a auditar despu\'es de abril. Algunos
conceptos pueden haber quedado desactualizados respecto de decisiones
posteriores ---por ejemplo, todo lo que se escribi\'o suponiendo compensaci\'on
por cancelaci\'on polo-cero, enfoque que se abandon\'o (Etapa~3).}

\section{Etapa 3 --- Electr\'onica anal\'ogica y modelo del transductor}
\label{sec:et3}
\periodo{17 de abril --- 20 de mayo de 2026 \quad$\cdot$\quad 9 d\'ias de trabajo registrados}

La etapa m\'as densa en decisiones de dise\~no. Arranca con el modelo del
ge\'ofono y termina con la cadena anal\'ogica que sigue en uso. Contiene los
\textbf{dos pivotes de arquitectura} del proyecto.

\subsection{El problema de partida}

El SM-24 tiene frecuencia natural $f_n = \SI{10}{\hertz}$ y amortiguamiento
$\zeta_0 = 0{,}25$. En la forma en aceleraci\'on su respuesta es

\begin{equation}
H_a(s) = \frac{-G_0\,s}{s^2 + 2\zeta_0\omega_n s + \omega_n^2},
\qquad \omega_n = 2\pi\cdot\SI{10}{\hertz}.
\label{eq:hgeo}
\end{equation}

Por debajo de $f_n$ la respuesta cae a \SI{20}{\decibel} por d\'ecada. El
problema es que la informaci\'on de las capas profundas del suelo vive
justamente ah\'i: mayor profundidad de investigaci\'on significa mayor longitud
de onda, y mayor longitud de onda significa menor frecuencia. Con el
transductor tal cual, la profundidad alcanzable est\'a acotada por el
transductor, no por el suelo ni por la electr\'onica.

\subsection{Primer intento: filtro RC pasivo de entrada (17 de abril)}

El primer ataque fue un filtro RC pasivo antes del amplificador, dise\~nado por
optimizaci\'on num\'erica (\texttt{FiltroPasivoEntrada.m}): se plante\'o la
funci\'on de transferencia conjunta ge\'ofono\,$+$\,RC ---que resulta de tercer
orden--- y se buscaron $R$ y $C$ que aplanaran la respuesta. Se explor\'o en
paralelo una alternativa por \textbf{observador de estado} en Simulink
(\texttt{SM24\_model\_Observer1.slx}).

Ambas cosas se abandonaron el 18 de abril, menos de 28 horas despu\'es.

\descartado{\textbf{Filtro RC pasivo de entrada.} Motivos registrados en la
bit\'acora: \textbf{(1)} la frecuencia de corte queda a merced de la tolerancia
de los componentes ---con \SIrange{5}{10}{\percent} t\'ipico, $f_c$ se corre lo
suficiente como para amplificar de m\'as la resonancia del ge\'ofono en vez de
compensarla---; \textbf{(2)} deriva t\'ermica de los condensadores en trabajo de
campo entre \SI{0}{\celsius} y \SI{40}{\celsius}; \textbf{(3)} \emph{fase no
lineal}, que para MASW es descalificante, porque el m\'etodo mide justamente
velocidad de fase y un error de fase se traduce directo en error de velocidad;
y \textbf{(4)} el PSoC ya trae DSP, as\'i que el filtrado sale gratis en
componentes.}

\pendiente{La bit\'acora del 18/04 presenta estos motivos como
\emph{reconstrucci\'on} a partir del \emph{diff}, no como testimonio directo
---dice literalmente ``las razones probables''--- y en un pasaje habla de un
ge\'ofono de \SI{4.5}{\hertz}, que no es el SM-24. El motivo (3), fase lineal,
es el t\'ecnicamente decisivo y el que sostiene el resto del dise\~no.
\textbf{Confirmar cu\'al fue el motivo real.}}

\subsection{Lo que reemplaz\'o al RC: FIR con decimaci\'on en cascada}

En lugar del RC se dise\~n\'o una cadena de tres etapas FIR de 32 coeficientes
(\texttt{fir1(31,\,0.15)}, ventana de Hamming), cada una seguida de decimaci\'on
$\times 4$ ---decimaci\'on total $\times 64$.

\begin{center}\small
\begin{tabularx}{\textwidth}{@{}l X@{}}
\toprule
\textbf{Propiedad} & \textbf{Por qu\'e importa ac\'a} \\
\midrule
Fase lineal exacta & Todas las frecuencias se retardan lo mismo: no contamina la curva de dispersi\'on. Es la raz\'on principal \\
Estabilidad incondicional & Solo ceros, sin polos \\
Coeficientes exactos & Sin tolerancia, sin deriva t\'ermica, sin envejecimiento \\
Reconfigurable & Cambiar $f_c$ es cambiar una l\'inea, no un componente \\
Costo computacional & Decimar en cascada cuesta $\approx 32(1+\tfrac14+\tfrac1{16})\approx 42$ MAC/muestra equivalentes, contra los \textasciitilde{}2000 coeficientes que exigir\'ia un \'unico filtro de $\times64$ \\
\bottomrule
\end{tabularx}
\end{center}

\seleccionado{\textbf{El principio qued\'o y estructura todo el sistema:} lo que
puede resolverse en digital con fase controlada se resuelve en digital; lo
anal\'ogico se reserva para lo que \emph{solo} puede hacerse antes de
cuantizar ---acondicionar el nivel y extender la banda del transductor.}

\subsection{Segundo intento: compensaci\'on activa}

Con el RC descartado, el problema de la banda baja segu\'ia intacto: ning\'un
filtro digital puede recuperar informaci\'on que el transductor atenu\'o por
debajo del piso de ruido \emph{antes} del ADC. La extensi\'on de banda tiene que
ocurrir en anal\'ogico.

Entre el 22 y el 30 de abril se probaron tres l\'ineas:

\begin{center}\small
\begin{tabularx}{\textwidth}{@{}l X l@{}}
\toprule
\textbf{Script} & \textbf{Enfoque} & \textbf{Destino} \\
\midrule
\texttt{compensadorKai.m} & Compensador anal\'itico de forma cerrada & Explorado \\
\texttt{compensadorPazos.m} & B\'usqueda exhaustiva sobre el espacio de componentes & Explorado \\
\texttt{Compensador.m} & Filtro activo VCVS / Sallen-Key & \textbf{Descartado 30/04} \\
\bottomrule
\end{tabularx}
\end{center}

La forma objetivo qued\'o fijada como

\begin{equation}
H_3(s) = \frac{2(\zeta_1-\zeta_0)\,\omega_n\,s}
              {s^2 + 2\zeta_1\omega_n s + \omega_n^2},
\label{eq:objetivo}
\end{equation}

con $\zeta_0=0{,}25$ ---el del ge\'ofono, que el numerador debe reproducir para
cancelar sus polos--- y $\zeta_1$ el amortiguamiento \emph{nuevo}, mucho mayor,
que fija el denominador de la etapa pasa-banda. En el script de abril se us\'o
$\zeta_1=100$ como valor de trabajo; la caracterizaci\'on medida de julio
(Etapa~9) da $\zeta_1 \approx 937$ para el circuito calibrado.

Esta es exactamente la estrategia de Ma \emph{et al.} (2023): \textbf{no se
cancelan polos ni ceros}, se sube el amortiguamiento manteniendo $\omega_n$.
La red completa resulta $1 - H_3(s)$ ---pasa-todo menos pasa-banda---, y \emph{eso
explica por qu\'e la cadena anal\'ogica del proyecto tiene una etapa pasa-banda
seguida de una suma}: el pasa-banda y la resta \textbf{son} la compensaci\'on.

\begin{figure}[!ht]
\centering
\includegraphics[width=0.86\textwidth]{extension_banda.png}
\caption{Qu\'e hace la compensaci\'on. Al elevar $\zeta_1$, los polos complejos
del transductor se separan en dos reales y el inferior queda en
$\omega_n/(2\zeta_1)$. El corte de \SI{-3}{\decibel} pasa de
\SI{8.0}{\hertz} sin compensar a \SI{0.51}{\hertz} con el $\zeta_1=10$ de Ma
\emph{et al.}\ y a \SI{0.010}{\hertz} con el $\zeta_1\approx 937$ que resulta
del denominador BP medido en la tanda calibrada del 21/07. Calculado a partir
de la Ec.~\eqref{eq:hgeo} y $1-H_3$ con $f_n = \SI{10.2046}{\hertz}$.
\textbf{Esta es la respuesta \emph{de dise\~no}}; la muesca realmente medida
resulta \SI{26.3}{\decibel} menos profunda de lo previsto (Etapa~9).}
\label{fig:extension}
\end{figure}

\subsection{El pivote de topolog\'ia: Sallen-Key $\to$ MFB (30 de abril)}

El 30 de abril se escribi\'o \texttt{ComparacionFiltros.m}, que superpone los
Bode de tres funciones: la MFB propuesta, la Sallen-Key vigente y la respuesta
objetivo de la Ec.~\eqref{eq:objetivo}. Ese gr\'afico fue la herramienta de
decisi\'on.

\begin{figure}[!ht]
\centering
\includegraphics[width=0.9\textwidth]{comparacion_topologias.png}
\caption{Reproducci\'on del \texttt{ComparacionFiltros.m} original, con los
mismos valores de componente. Arriba, las tres respuestas; abajo, la
desviaci\'on contra el objetivo. En \SIrange{0.1}{200}{\hertz} la MFB se aparta
\SI{0.19}{\decibel} como m\'aximo y la Sallen-Key \SI{0.26}{\decibel}: en
magnitud son intercambiables. \textbf{El pivote no se decidi\'o por respuesta.}}
\label{fig:topologias}
\end{figure}

\seleccionado{\textbf{Topolog\'ia MFB (\emph{Multiple Feedback}).} Criterio, en
orden de peso:
\begin{enumerate}[itemsep=0.1em,topsep=0.2em]
\item \textbf{Menor sensibilidad de $Q$ y $\omega_0$ a las tolerancias.} Es lo
decisivo: se usan componentes de tolerancia floja y hay que replicar varios
canales apareados. La Sallen-Key realimenta en positivo y por lo tanto
\emph{amplifica} el error de componentes; la MFB usa solo realimentaci\'on
negativa.
\item Mejor atenuaci\'on en alta frecuencia.
\item La baja impedancia de entrada de la MFB no penaliza, porque hay etapas
de ganancia antes de ella.
\item La inversi\'on de fase se revierte en digital sin costo.
\item El producto ganancia-ancho de banda de los operacionales es enorme
frente a las frecuencias Rayleigh: no es limitante.
\end{enumerate}
Dicho de otro modo: \emph{las desventajas cl\'asicas de la MFB son inocuas en
esta banda, y su ventaja ataca el problema real}.}

\descartado{\textbf{Topolog\'ia Sallen-Key / VCVS.} No por respuesta ---la
Fig.~\ref{fig:topologias} muestra que cumpl\'ia--- sino por sensibilidad a
tolerancias, riesgo de oscilaci\'on con ganancia mayor que uno, e impedancia de
salida real no nula del operacional.}

\begin{figure}[!ht]
\centering
\includegraphics[width=\textwidth]{AnalogPathFinal_trim.png}
\caption{La cadena anal\'ogica que result\'o de esta etapa y sigue vigente:
PGA $\to$ pasa-banda (MFB) $\to$ sumador en corriente $\to$ pasa-bajos, con
referencia ajustable por etapa. El pasa-banda y la suma \emph{son} la red de
compensaci\'on $1-H_3$: no son dos bloques que casualmente est\'an en serie.}
\label{fig:cadena}
\end{figure}

El mismo d\'ia entr\'o la otra decisi\'on de la etapa: la \textbf{suma en
corriente} mediante amplificador de transimpedancia. M\'ultiples se\~nales
entran por resistencias a un nodo de tierra virtual y el TIA convierte la suma
de corrientes en tensi\'on. Es un sumador lineal que no carga las fuentes y que
escala agregando una resistencia por entrada ---esa es la etapa SUM de la
cadena definitiva.

\subsection{Medici\'on: MFB contra Sallen-Key en banco}

La decisi\'on por an\'alisis se contrast\'o despu\'es con datos. Se capturaron
conjuntos con las dos topolog\'ias (\texttt{MFB\_DigitalFilterLPF.mat} y
\texttt{Sallen\_Key.mat}) y se construy\'o \texttt{analisis\_espectral.m} (264
l\'ineas) como banco de comparaci\'on:

\begin{itemize}[itemsep=0.2em]
\item \textbf{PSD} cruda y filtrada, cada medici\'on en gris y la media
$\pm\sigma$ en color, con las bandas de inter\'es resaltadas;
\item \textbf{SNR por sesi\'on}, para identificar capturas at\'ipicas y
descartarlas con criterio en vez de a ojo;
\item \textbf{espectrogramas}, que hacen visible el evento transitorio ---un
auto pasando aparece como una mancha de energ\'ia en \SIrange{2}{20}{\hertz}
durante \SIrange{1}{3}{\second}.
\end{itemize}

\dato{\textbf{Vol\'umen experimental de la etapa:} 18 sesiones de captura el 28
de abril, 41 sesiones entre el 2 y el 4 de mayo. El patr\'on de tiempos del 4 de
mayo alterna micro-sesiones de \SIrange{3}{16}{\second} ---cambiar un
componente, reconectar, verificar--- con sesiones largas de
\SIrange{30}{60}{\minute} de captura estable.}

\subsection{Primeras se\~nales reales y primera salida a exteriores}

El 20 de abril lleg\'o la primera se\~nal real del SM-24 a la PC, v\'ia
\texttt{geophone\_scope.m} (657 l\'ineas), un osciloscopio virtual en MATLAB
sobre puerto serie. El 29 de abril el ge\'ofono sali\'o por primera vez al
exterior ---las capturas \texttt{mesitaEstacionamiento*.mat}, unos 13 segundos
de se\~nal cada una, en el estacionamiento de la Facultad: un sitio con fuentes
de vibraci\'on abundantes y predecibles (autos, personas caminando) a pocos
metros del laboratorio.

\subsection{La cadena de firmware toma forma}

Para el 29 de abril el PSoC ya ten\'ia el \emph{pipeline} completo
ADC\,$\to$\,DFB\,$\to$\,FIR\,$\to$\,UART, con:

\begin{itemize}[itemsep=0.2em]
\item cuatro configuraciones de ADC seleccionables;
\item el \textbf{DFB} (filtro digital en hardware del PSoC) filtrando la salida
del $\Delta\Sigma$ antes del FIR por software;
\item un \textbf{notch FIR de \SI{50}{\hertz}} contra la interferencia de red
(\texttt{notch.m});
\item protocolo binario propio hacia la PC: \texttt{0xAA} muestra de 12 bytes,
\texttt{0xCC} confirmaci\'on de estado, \texttt{0xDD} evento de \emph{debug},
y comandos \texttt{[0xBB][CMD][LEN][payload]} en sentido inverso;
\item \emph{heartbeat} autom\'atico cada \textasciitilde{}\SI{2}{\second} para
detectar cuelgues.
\end{itemize}

\descartado{\textbf{El notch de \SI{50}{\hertz} termin\'o retirado} (decisi\'on
posterior, Etapa~9). La interferencia de red es determinista y de forma
conocida, as\'i que conviene \emph{estimarla y restarla} en post-proceso por
ajuste no lineal de m\'inimos cuadrados: eso no introduce retardo de grupo, es
reversible y ---lo importante--- no borra se\~nal s\'ismica \'util, que en
\SI{50}{\hertz} todav\'ia la hay.}

\subsection{Calibraci\'on manual del \emph{offset} (10 de mayo)}

``Calibraci\'on a manita ya funciona''. El \emph{offset} de continua de la
cadena diferencial se cancel\'o desplazando a mano las referencias de los VDAC,
con un protocolo de configuraci\'on PSoC$\leftrightarrow$MATLAB
(\emph{request} \texttt{0xA5}) para leer y fijar par\'ametros en caliente.

\seleccionado{Este es el germen directo de la auto-calibraci\'on de la
Etapa~6: primero se demostr\'o que ajustando las referencias se
cancelaba el \emph{offset}, y despu\'es se cerr\'o el lazo para que lo hiciera el
firmware solo.}

\subsection{Cierre del compensador: del ideal al circuito real (17--20 de mayo)}

Hasta ac\'a el compensador exist\'ia como funci\'on de transferencia. Faltaba
convertirlo en un circuito con componentes comprables y operacionales reales,
y eso ocup\'o tres d\'ias duros ---el mensaje de \emph{commit} del 18 de mayo
fue \emph{``ayudame dios anal\'ogico quiero sucumbir al diablo digital''}.

\textbf{El operacional real entra al modelo.} Se dej\'o de suponer ganancia
infinita: se model\'o el amplificador con $A_0 = \SI{90}{\decibel}$ y polo
dominante en \SI{8}{\mega\hertz}, y se reverific\'o la respuesta del
compensador con esa no idealidad
(\texttt{graficar\_compensador\_csv.py}). Este es el paso que suele delatar un
dise\~no que ``andaba en el papel''.

\textbf{B\'usqueda de componentes.} Se escribieron cuatro variantes de
optimizaci\'on en paralelo, atacando el mismo problema desde \'angulos
distintos. La que qued\'o plantea el espacio de b\'usqueda como continuo en tres
dimensiones $[\log_{10}R_1, \log_{10}C_1, \log_{10}C_2]$ y luego:

\begin{enumerate}[itemsep=0.15em]
\item determina $R_2$ imponiendo la condici\'on de $\omega_0$ exacta ---o sea,
elimina una dimensi\'on por restricci\'on en vez de por b\'usqueda;
\item resuelve las razones del sumador anal\'iticamente por m\'inimos cuadrados;
\item reci\'en al final hace \emph{snap} a valores comerciales y resuelve lo que
no exista con redes en T.
\end{enumerate}

\seleccionado{El orden importa y es la lecci\'on transferible: \emph{optimizar
en el continuo, restringir por f\'isica, discretizar al final}. Discretizar
antes de optimizar deja el \'optimo fuera de la grilla.}

\textbf{La cadena entra al PSoC (20 de mayo).} Con el proyecto
\texttt{DiferencialToSingleEnded} se implement\'o por primera vez la cadena
anal\'ogica completa dentro del PSoC ---el \texttt{main.c} que arranca todos los
bloques anal\'ogicos en la secuencia correcta. A la \'una de la ma\~nana, con
osciloscopio conectado: \emph{``Datos y m\'as datos, empieza la locura''}.

\dato{Se agreg\'o adem\'as un AMUX para poder \emph{medir el modo com\'un} de la
entrada diferencial, y el protocolo se comprimi\'o a 4 bytes multiplexando tipo
y canal en el segundo byte.}

\subsection{Balance de la etapa}

\begin{center}\small
\begin{tabularx}{\textwidth}{@{}l X@{}}
\toprule
\textbf{Se qued\'o} & Compensaci\'on activa por elevaci\'on de $\zeta$ (no cancelaci\'on polo-cero); topolog\'ia MFB; suma en corriente por TIA; filtrado de fase lineal en digital; DFB en hardware; protocolo binario propio; ajuste de \emph{offset} por VDAC \\
\addlinespace[3pt]
\textbf{Se descart\'o} & Filtro RC pasivo de entrada; observador de estado; topolog\'ia Sallen-Key/VCVS; notch de \SI{50}{\hertz} (m\'as tarde) \\
\addlinespace[3pt]
\textbf{Qued\'o abierto} & El \emph{offset} se cancelaba a mano y hab\'ia que repetirlo por placa y por ganancia: no escalaba a un arreglo de varios nodos \\
\bottomrule
\end{tabularx}
\end{center}

\section{Etapa 4 --- La red inal\'ambrica ESP-NOW}
\label{sec:et4}
\periodo{24 de mayo --- 3 de junio de 2026 \quad$\cdot$\quad 7 d\'ias de trabajo, 18 \emph{commits}}

\subsection{El problema que resuelve}

Un sism\'ografo multicanal convencional cablea todos los ge\'ofonos a un
digitalizador \'unico, y entonces la sincronizaci\'on entre canales es gratis:
hay un solo reloj. Al pasar a nodos inal\'ambricos independientes, cada nodo
tiene su propio reloj y su propio arranque, y \emph{la sincronizaci\'on pasa a
ser el problema central del sistema}.

Y no es un problema menor de comodidad. En MASW, un desfasaje entre canales se
traduce directamente en un error de velocidad de fase, que se propaga a la
curva de dispersi\'on y de ah\'i al perfil $V_S$. Un arreglo desincronizado no
da un resultado ruidoso: da un resultado \emph{equivocado y cre\'ible}.

\subsection{Arquitectura adoptada}

\begin{center}\small
\begin{tabularx}{\textwidth}{@{}l X@{}}
\toprule
\textbf{Elemento} & \textbf{Rol} \\
\midrule
Maestro ESP32 & Coordina el arreglo, sincroniza el disparo, concentra los vol\'umenes de datos y expone la interfaz hacia la PC \\
Esclavos ESP32 & Uno por nodo. Comandan al PSoC local, reciben las muestras y las retienen hasta el volcado \\
ESP-NOW & Enlace de control y sincronizaci\'on: sin asociaci\'on previa, latencia baja y predecible ---lo que hace falta para disparar, no para transferir \\
PSoC & Adquisici\'on real. El ESP no toca la se\~nal: pide, espera y transporta \\
\bottomrule
\end{tabularx}
\end{center}

\descartado{\textbf{Enlace SPI PSoC\,$\leftrightarrow$\,ESP.} La primera
versi\'on (24 de mayo) usaba SPI con el PSoC como esclavo. Se reemplaz\'o por
UART en el refactor del 31 de mayo. En la placa de agosto la subida vuelve a
cambiar, ahora a I\textsuperscript{2}C (Etapa~11): es la interfaz que
m\'as veces se revis\'o del proyecto.}

\subsection{La secuencia de sincronizaci\'on: c\'omo se lleg\'o a PRESTART}

Vale la pena registrar la evoluci\'on completa, porque es el mejor ejemplo de
c\'omo el sistema aprendi\'o a no confiar en el tiempo.

\textbf{Intento 1 --- espera fija.} \verb|ARM| $\to$ esperar $N$ ms $\to$
\verb|START|. Falla cuando ESP-NOW tarda m\'as que $N$: el esclavo recibe
\verb|START| antes de estar en \verb|ARMED| y directamente no arranca. El
s\'intoma es traicionero, porque el resto del arreglo s\'i arranca.

\textbf{Intento 2 --- \emph{probe}/ACK con medici\'on de RTT} (25 de mayo). Ya
no se supone la latencia: se mide. Mejora, pero sigue habiendo ventana entre
la confirmaci\'on y el disparo.

\textbf{Soluci\'on --- PRESTART con HOT\_WAIT} (26 de mayo). Se agrega un estado
intermedio en el que el maestro \emph{verifica} que todos los esclavos est\'an
calientes antes de emitir el disparo real:

\begin{center}\small
\texttt{IDLE $\to$ ARMING $\to$ ARMED $\to$ \textbf{PRESTART} $\to$ RUNNING $\to$ STOPPING}
\end{center}

\begin{enumerate}[itemsep=0.15em]
\item El maestro difunde \verb|CMD_ARM|.
\item Espera los \verb|ARM_ACK| con \emph{timeout}.
\item Entra en PRESTART y consulta a cada esclavo su estado (HOT\_WAIT:
consulta cada \SI{20}{\milli\second}, \emph{timeout} total
\SI{120}{\milli\second}, 3 reintentos, todo configurable).
\item Solo con todos confirmados emite el \verb|START| real.
\end{enumerate}

\seleccionado{\textbf{Y todo el mecanismo es no bloqueante.} Este es el detalle
que lo hizo funcionar: los pulsos de \emph{debug} usaban
\texttt{delayMicroseconds()}, lo que congelaba el CPU \SI{1}{\milli\second} y
hac\'ia perder callbacks de ESP-NOW ---es decir, \emph{el propio instrumento de
medici\'on causaba la p\'erdida del \emph{ACK} que se estaba tratando de
diagnosticar}. Se reescribi\'o con marcas de tiempo y verificaci\'on en el
\emph{loop}. Cambio chico en c\'odigo, decisivo en comportamiento.}

\subsection{C\'omo se midi\'o la sincronizaci\'on}

No se dio por buena: se instrument\'o. Se agreg\'o el comando
\verb|CMD_SCOPE_START| (\texttt{0x60}), un disparo \emph{falso} que hace que
todos los nodos emitan un pulso GPIO de \SI{1}{\milli\second} sin muestrear ni
responder. Con el osciloscopio en los pines de \emph{sync} de los tres nodos
(maestro GPIO\,15, esclavo\,1 GPIO\,32, esclavo\,2 GPIO\,22), la separaci\'on
entre pulsos \emph{es} la latencia de sincronizaci\'on del sistema.

Y para no medir una sola vez, el estado \verb|SCOPE_MULTI| repite el disparo
128 veces con \SI{1}{\second} de separaci\'on, de modo que el osciloscopio
acumule estad\'istica de \emph{jitter} sin intervenci\'on manual.

\dato{Objetivo de dise\~no fijado en esta etapa: todos los pulsos dentro de
\SI{1}{\milli\second} entre s\'i.}

\pendiente{El \emph{jitter} medido sobre el arreglo completo y desplegado en
campo sigue sin volcarse al documento. Ver \S\ref{sec:pendientes}.}

\subsection{Hitos de funcionamiento}

\begin{center}\small
\begin{tabularx}{\textwidth}{@{}l l X@{}}
\toprule
\textbf{Fecha} & \textbf{Mensaje} & \textbf{Qu\'e significa} \\
\midrule
26 may & \emph{``Parece que funciona la putaaaa''} & Primera transmisi\'on extremo a extremo: PSoC $\to$ ESP esclavo $\to$ ESP maestro $\to$ MATLAB \\
3 jun & \emph{``Locura parece que ya funciona super bien''} & Sistema estable tras corregir el \emph{beacon} a \SI{10}{\hertz}, \emph{store-and-forward} en \texttt{VER} y UART a 115\,200~baudios \\
\bottomrule
\end{tabularx}
\end{center}

\subsection{En paralelo: acondicionamiento del martillo (28--29 de mayo)}

La fuente de impacto necesita su propio canal, porque marca el instante cero
de la captura. Se cre\'o el proyecto PSoC \texttt{AcondicionamientoMartillo}
para el aceler\'ometro de impacto \textbf{PCB Piezotronics 086D20}, con cadena
propia: \texttt{LPF\_ref} $+$ \texttt{LPF\_ADC} $+$ referencia de
\SI{1.024}{\volt}.

Que sea un proyecto aparte no es casual: el martillo y el ge\'ofono tienen
rangos de se\~nal y anchos de banda muy distintos, y compartir cadena habr\'ia
obligado a comprometer los dos.

\subsection{Refactor de cierre (31 de mayo, PR \#2)}

El PR \#2 consolid\'o la etapa:

\begin{itemize}[itemsep=0.2em]
\item enlace PSoC\,$\leftrightarrow$\,ESP migrado de SPI a \textbf{UART};
\item contador de muestras $N$ ampliado a 16 bits;
\item m\'aquina de estados expl\'icita en el PSoC;
\item \textbf{arranque del PSoC solo por flanco f\'isico} ---no por comando---,
que es lo que garantiza que el instante cero sea el mismo en todos los nodos
independientemente de la latencia del enlace;
\item \texttt{debug\_log.h} unificado y READMEs por proyecto.
\end{itemize}

\seleccionado{\textbf{El principio que sobrevivi\'o a todo:} el disparo viaja
por software, pero el instante cero lo define un \emph{flanco f\'isico}. Todos
los rediseños posteriores ---incluida la placa de agosto, que cambia el enlace
de subida a I\textsuperscript{2}C--- mantuvieron esa separaci\'on.}

\subsection{Balance de la etapa}

\begin{center}\small
\begin{tabularx}{\textwidth}{@{}l X@{}}
\toprule
\textbf{Se qued\'o} & ESP-NOW para control y sincronizaci\'on; PRESTART/HOT\_WAIT; disparo por flanco f\'isico; instrumentaci\'on de \emph{jitter} por pulso GPIO; \texttt{SCOPE\_MULTI} para estad\'istica autom\'atica; UART hacia el PSoC; canal dedicado de martillo \\
\addlinespace[3pt]
\textbf{Se descart\'o} & Espera fija tras ARM; enlace SPI PSoC$\leftrightarrow$ESP; pulsos de \emph{debug} bloqueantes \\
\addlinespace[3pt]
\textbf{Qued\'o abierto} & La captura segu\'ia gobernada por c\'odigo C en el PSoC, con el \emph{jitter} atado a las interrupciones. Ese es el punto de partida de la Etapa~5 \\
\bottomrule
\end{tabularx}
\end{center}

\section{Etapa 5 --- Consolidaci\'on del nodo e interfaz de campo}
\label{sec:et5}
\periodo{4 --- 10 de junio de 2026 \quad$\cdot$\quad 6 d\'ias, 9 \emph{commits}}

Etapa corta y de mucha consecuencia: el sistema deja de ser un banco de
laboratorio y empieza a parecerse a un instrumento que se puede llevar al
campo.

\begin{figure}[ht]
\centering
\includegraphics[width=\textwidth]{DigitalPathFinal_trim.png}
\caption{Camino digital consolidado en PSoC Creator: ADC, DFB, DMA y control en
hardware. Documenta la arquitectura implementada, no el rechazo anti-\emph{alias} de
cada factor de decimaci\'on configurable.}
\end{figure}

\subsection{Un solo firmware para dos nodos distintos}

El 4 de junio se unific\'o el firmware. Hasta ese momento hab\'ia dos ramas de
c\'odigo: el nodo ge\'ofono y el nodo martillo. Pasaron a compartir un \'unico
\texttt{main.c} que \textbf{detecta en tiempo de compilaci\'on qu\'e hardware
tiene debajo}, inspeccionando qu\'e bloques declara el TopDesign del PSoC, y se
compila como GEO o como HAMMER en consecuencia.

\seleccionado{\textbf{Es la decisi\'on de arquitectura de software m\'as
rentable del proyecto.} Elimina la bifurcaci\'on del firmware embebido: el
protocolo, el \emph{store-and-forward}, la m\'aquina de estados y el manejo de
errores tienen \emph{una sola} implementaci\'on. Un arreglo con nodos que
divergen en firmware es un arreglo que se desincroniza de maneras dif\'iciles
de diagnosticar; ac\'a eso directamente no puede pasar. Lo \'unico que difiere
entre nodos es la cadena anal\'ogica, que es donde \emph{debe} diferir.}

\subsection{La PC deja de ser necesaria (8 de junio)}

El maestro ESP32 pas\'o a servir una aplicaci\'on web completa sobre su propio
punto de acceso WiFi, con relay WebSocket para telemetr\'ia en vivo.

Esto cambia el modo de operaci\'on en campo: ya no hace falta llevar una
notebook con MATLAB y un cable USB. Se levanta el AP del maestro, se conecta
el tel\'efono y se opera desde ah\'i. Para un instrumento que se despliega en
terreno, esa diferencia es pr\'actica, no cosm\'etica.

\subsection{El notch de \SI{50}{\hertz} se retira (7 y 9 de junio)}

Aqu\'i ocurre el descarte que se hab\'ia anticipado en la Etapa~3,
y ocurre en dos pasos:

\descartado{\textbf{Paso 1 (7 de junio): fuera el LMS adaptativo.} Se hab\'ia
probado cancelar la interferencia de red con un filtro adaptativo LMS. Se
reemplaz\'o por estimaci\'on directa por m\'inimos cuadrados. Motivo: el LMS
\emph{converge} ---tiene transitorio de adaptaci\'on--- y en una captura
s\'ismica de pocos segundos el transitorio ocupa una fracci\'on no despreciable
del registro; adem\'as su comportamiento depende de la se\~nal, lo que
introduce una no linealidad efectiva que arruina la comparabilidad entre
canales.}

\descartado{\textbf{Paso 2 (9 de junio): fuera el notch arm\'onico del cliente
web.} La \SI{50}{\hertz} cae \emph{dentro} de la banda \'util del MASW, as\'i que
un notch borra se\~nal junto con la interferencia, e introduce retardo de grupo
justo donde se mide fase.}

\seleccionado{\textbf{Lo que qued\'o:} una b\'usqueda acotada de la frecuencia
de red, seguida por una regresi\'on lineal seno--coseno de la fundamental y sus
arm\'onicos sobre el registro completo. La componente estimada se resta en
post-proceso: no agrega retardo de grupo y el dato crudo se conserva. No es
neutro frente a una componente s\'ismica exactamente coherente con las bases
ajustadas, que tambi\'en puede atenuarse; esa limitaci\'on debe quedar visible.}

En el mismo movimiento se subi\'o la tasa del ADC a \SI{3}{\kilo\hertz}.

\subsection{El \emph{handoff} de \emph{jitter} (10 de junio)}

Documento de an\'alisis de robustez del sistema de sincronizaci\'on. Contiene el
hallazgo m\'as fino de la etapa:

\dato{\textbf{Bug ACK-antes-de-GPIO.} El esclavo confirmaba la recepci\'on del
disparo \emph{antes} de haber generado el flanco f\'isico que realmente arranca
al PSoC. El maestro entonces cre\'ia sincronizado un arreglo que a\'un no lo
estaba, y el error no aparec\'ia como falla sino como \emph{jitter} ---es decir,
como una peque\~na incertidumbre en el instante cero de cada canal, exactamente
la cantidad que MASW convierte en error de velocidad de fase.}

La correcci\'on vigente es m\'as estricta que ``mover el ACK'': ante un
\texttt{CMD\_START} real, el esclavo eleva primero
\texttt{SYNC\_TO\_PSOC\_PIN} y, desde ese instante, permanece en silencio de
radio hasta completar la captura. Los \texttt{START\_ACK} inmediatos quedan
para sondeos o rechazos previos al muestreo, no para confirmar una captura en
curso.

El documento adem\'as plantea el \textbf{presupuesto de \emph{jitter} entre
nodos}, que es la forma correcta de encarar el problema: en vez de perseguir
``que sincronice bien'', se enumeran las contribuciones (latencia ESP-NOW,
resoluci\'on del disparo, deriva de reloj, retardo del flanco) y se acota cada
una.

\pendiente{El presupuesto est\'a planteado te\'oricamente y hubo una observaci\'on
preliminar inferior a \SI{400}{\micro\second} sobre pocos nodos a unos
\SI{5}{\meter}. No es una especificaci\'on: tras retirar el cable de
sincronismo falta medir la ruta completa callback de radio$\to$flanco
local$\to$inicio del PSoC sobre el arreglo desplegado.}

\subsection{Balance de la etapa}

\begin{center}\small
\begin{tabularx}{\textwidth}{@{}l X@{}}
\toprule
\textbf{Se qued\'o} & Firmware \'unico con detecci\'on de hardware en compilaci\'on; interfaz web embebida en el maestro (operaci\'on sin PC); cancelaci\'on de red por ajuste de m\'inimos cuadrados en post-proceso; presupuesto de \emph{jitter} como m\'etodo \\
\addlinespace[3pt]
\textbf{Se descart\'o} & Firmware bifurcado GEO/HAMMER; filtro adaptativo LMS; notch arm\'onico en el cliente; ACK emitido antes del flanco f\'isico \\
\bottomrule
\end{tabularx}
\end{center}

\section{Etapa 6 --- Auto-calibraci\'on, DMA y \texttt{superMaquina}}
\label{sec:et6}
\periodo{14 de junio --- 2 de julio de 2026 \quad$\cdot$\quad 11 d\'ias, 17 \emph{commits}}

La etapa que produjo el resultado publicable del proyecto. Tiene tres l\'ineas
que avanzan entrelazadas: cerrar el lazo de calibraci\'on del \emph{front-end},
sacar la adquisici\'on de la CPU, y finalmente sacar la \emph{decisi\'on} de
adquirir tambi\'en de la CPU.

\subsection{Por qu\'e hace falta calibrar}

La cadena anal\'ogica tiene cuatro etapas en cascada ---PGA, pasa-banda (BP),
sumador (SUM) y pasa-bajos (LP)--- y cada una aporta su propio \emph{offset} de
continua: tensiones de \emph{offset} de los operacionales, corrientes de
polarizaci\'on, desapareamiento de resistencias. El problema es que
\emph{el \emph{offset} se amplifica junto con la se\~nal}, y la compensaci\'on
del ge\'ofono exige ganancia alta ---proporcional a $(\zeta_1-\zeta_0)$, con
$\zeta_0 = 0{,}25$ y $\zeta_1$ de varios cientos: tres \'ordenes de magnitud.

\dato{Dimensi\'on del problema, medida: el \emph{offset} con referencia fija a
la salida de la etapa sumadora lleg\'o a \textbf{\SI{21.8}{\percent}} del fondo
de escala, y a la salida del LP ---la que alimenta al ADC--- hasta
\textbf{\SI{45.8}{\percent}}. Es decir, casi la mitad del rango de entrada
consumido \emph{antes de aplicar se\~nal s\'ismica alguna}.}

\subsection{Los tres algoritmos que se probaron}

\begin{center}\small
\begin{tabularx}{\textwidth}{@{}l X l@{}}
\toprule
\textbf{Enfoque} & \textbf{Qu\'e pas\'o} & \textbf{Destino} \\
\midrule
Bisecci\'on binaria sobre c\'odigos VDAC
& Primer intento (14 jun), solo para GEO. Funciona pero es ciego: no aprovecha que la planta es conocida, y cada etapa necesitaba su propia versi\'on
& \textbf{Descartado} \\
\addlinespace[2pt]
Servo lento
& Convergencia por integraci\'on pura
& \textbf{Descartado} \\
\addlinespace[2pt]
PID
& Se prob\'o y se abandon\'o. El t\'ermino derivativo amplifica el ruido de medici\'on del ADC sin aportar nada en una planta cuasi-est\'atica ---en la bit\'acora qued\'o descrito como \emph{``overfitting al ruido''}
& \textbf{Descartado el D} \\
\addlinespace[2pt]
\textbf{PI con \emph{feedforward} f\'isico}
& Un \'unico controlador para las cuatro etapas de GEO y las dos de HAMMER
& \textbf{Se qued\'o} \\
\bottomrule
\end{tabularx}
\end{center}

\seleccionado{\textbf{El 23 de junio los dos algoritmos separados (GEO y
HAMMER) convergen en uno solo: PI puro con \emph{feedforward}}, detector de
\emph{lock} portado de una simulaci\'on de Simulink de referencia. Las
implementaciones de bisecci\'on y servo no se borraron de inmediato: quedaron
compiladas fuera con \texttt{\#if 0} como referencia, y se eliminaron del
archivo activo reci\'en el 26 de junio, una vez validado el reemplazo. Ese
orden ---validar antes de borrar--- vale la pena registrarlo.}

\begin{figure}[!ht]
\centering
\includegraphics[width=0.94\textwidth]{psoc_foreground_calibration_flow.drawio_horizontal.png}
\caption{Flujo de la calibraci\'on \emph{foreground} en el PSoC: las etapas se
calibran en orden, cada una hasta entrar en su banda muerta, antes de pasar a
la siguiente. Al terminar, el resultado se persiste en EEPROM indexado por
c\'odigo de ganancia.}
\label{fig:calib-flow}
\end{figure}

\begin{figure}[!ht]
\centering
\includegraphics[width=0.80\textwidth]{calibration_fsm_simplified.png}
\caption{M\'aquina de estados de la calibraci\'on, versi\'on simplificada. La
versi\'on detallada est\'a en \texttt{docs/calibration\_fsm/}.}
\label{fig:calib-fsm}
\end{figure}

\subsection{Los detalles finos del controlador}

Tres afinaciones sucesivas, cada una corrigiendo un modo de falla real:

\begin{enumerate}[itemsep=0.3em]
\item \textbf{Escala en c\'odigos de DAC, no en volts} (24 jun). El actuador es
un DAC; expresar el lazo en su unidad natural evita un factor de conversi\'on
que dependía de la ganancia de etapa.

\item \textbf{Banda muerta proporcional a la ganancia de etapa} (24 jun). No
tiene sentido exigir la misma precisi\'on absoluta a una etapa de ganancia
unitaria que a una de ganancia 32: en la primera, insistir solo agrega ciclos
de conmutaci\'on del DAC sin mejorar nada aguas abajo.

\item \textbf{Anti-\emph{windup} direccional} (25 jun). Este es el m\'as sutil.
El integrador no se congela ante \emph{cualquier} saturaci\'on contra el riel
del DAC, sino \'unicamente cuando seguir integrando \emph{en ese sentido}
empeorar\'ia la saturaci\'on. Congelarlo en ambos sentidos deja al lazo incapaz
de salir del riel. Adem\'as hubo que manejar correctamente el \textbf{signo de
la ganancia} por etapa: \texttt{HAMMER\_PGA} es inversora y din\'amica,
\texttt{HAMMER\_LP} es fija y no inversora.
\end{enumerate}

\dato{Primeras convergencias documentadas en hardware real:
\texttt{HAMMER\_PGA ok=1 dac=65}, \texttt{HAMMER\_LP ok=1 dac=57} (23 jun); al
d\'ia siguiente \texttt{HAMMER\_LP} converge con \texttt{err\_mV=0} exacto
contra un objetivo de \SI{1024}{\milli\volt}. El 26 de junio el PI corre por
primera vez sobre las cuatro etapas de GEO.}

\subsection{Dos bugs que valen su propia entrada}

\dato{\textbf{El filtro cargaba los coeficientes equivocados} (24 jun). Se
descubri\'o que el filtro de adquisici\'on llevaba \emph{meses} cargando
coeficientes de pasa-banda donde deb\'ia haber un pasa-bajos. No se manifestaba
como falla evidente: el sistema funcionaba, simplemente filtraba otra cosa.}

\dato{\textbf{Un pin sin soldar} (26 jun). Un problema de calibraci\'on que se
persigui\'o como bug de firmware result\'o ser hardware. Registrado porque es la
clase de error que consume horas de depuraci\'on en el dominio equivocado.}

\subsection{Persistencia: la calibraci\'on sobrevive al reset}

El 17 de junio la calibraci\'on empez\'o a guardarse en EEPROM con CRC-16. El
1\textsuperscript{o} de julio se rehizo el esquema para guardar
\textbf{un \emph{slot} por c\'odigo de ganancia del PGA}, porque el punto de
operaci\'on \'optimo depende de la ganancia seleccionada.

\seleccionado{Consecuencia operativa: con calibraci\'on almacenada v\'alida, el
arranque solo ejecuta los intervalos de establecimiento por etapa
(2560 muestras, \textasciitilde{}\SI{1.0}{\second} a \SI{2604}{\hertz}) en vez
de la calibraci\'on completa. La diferencia entre esperar un segundo y esperar
un minuto es la diferencia entre un instrumento usable en campo y uno que no.}

\subsection{La adquisici\'on sale de la CPU}

Tres pasos en tres d\'ias:

\begin{center}\small
\begin{tabularx}{\textwidth}{@{}l l X@{}}
\toprule
\textbf{Fecha} & \textbf{Cambio} & \textbf{Consecuencia} \\
\midrule
17 jun & ISR cl\'asica $\to$ cadena DMA con \texttt{Filter} FIR de hardware &
El ARM deja de intervenir en el camino de la muestra. El FIR corre en el bloque dedicado, no en software \\
\addlinespace[2pt]
18 jun & Multiplexor DMA de cuatro canales (crudo/filtrado $\times$ simple/lote) gobernado por \texttt{Reg\_Select} &
Se cablea el promediado por hardware pero se lo deja inerte: la ruta existe antes de usarse \\
\addlinespace[2pt]
19 jun & Se activa el canal ``Lote'' &
La calibraci\'on deja de promediar muestra a muestra en el ISR y consume ventanas ya sumadas por hardware, con verificaci\'on \emph{en tiempo de compilaci\'on} de que el tama\~no del lote DMA coincide con la ventana de calibraci\'on \\
\bottomrule
\end{tabularx}
\end{center}

\subsection{\texttt{superMaquina}: la captura se vuelve hardware (2 de julio)}

El paso final. La l\'ogica de captura ---que hasta ese momento viv\'ia en C,
como ISRs sobre \texttt{Reg\_Select}--- se reescribi\'o como una \textbf{m\'aquina
de estados en Verilog dentro de un bloque UDB} del PSoC.

\texttt{superMaquina} decide, en hardware:

\begin{itemize}[itemsep=0.2em]
\item qu\'e DRQ de DMA se deja pasar;
\item cu\'ando arranca y cu\'ando termina una captura;
\item qu\'e interrupciones llegan efectivamente al ARM.
\end{itemize}

\seleccionado{\textbf{Por qu\'e importa tanto para MASW.} Mientras la captura la
gobierna software, el instante de inicio y el intervalo entre muestras heredan
el \emph{jitter} de las interrupciones ---latencia variable, prioridades,
secciones cr\'iticas. En hardware, el temporizado es determinista por
construcci\'on. Como el m\'etodo mide velocidad de fase, y la fase se mide contra
el tiempo, \emph{determinismo temporal es exactitud de medici\'on}.}

\begin{figure}[!ht]
\centering
\includegraphics[width=\textwidth]{DigitalPathFinal_trim.png}
\caption{Camino digital dentro del PSoC tras esta etapa: ADC $\Delta\Sigma$
$\to$ \texttt{Filter} FIR de hardware $\to$ cadena DMA con canales
crudo/filtrado y simple/lote, todo gobernado por \texttt{superMaquina}. El ARM
queda fuera del camino de la muestra.}
\label{fig:digital}
\end{figure}

En paralelo, tres \emph{timers} fijos del PSoC reemplazaron el \emph{tick} de
sistema de \SI{10}{\milli\second} (ping de arranque/reposo, parpadeo de LED,
\emph{watchdog} y progreso de calibraci\'on), y el ESP32 esclavo gan\'o un juego
completo de comandos por USB para operar el \emph{pipeline} paso a paso desde
la terminal, sin depender de la web: \texttt{probe}, \texttt{status},
\texttt{stream}, \texttt{debugpsoc}, \texttt{pre}, \texttt{sync},
\texttt{startnow}, \texttt{cap}, \texttt{clear}, \texttt{stop}.

\seleccionado{Ese juego de comandos result\'o ser la herramienta de diagn\'ostico
m\'as usada del proyecto: permite reproducir en aislamiento cualquier paso de
la secuencia de captura, que es exactamente lo que hizo falta para encontrar
los bugs cr\'iticos de la Etapa~8.}

\pendiente{Pol\'itica de eventos que qued\'o fijada y conviene no romper: el
\texttt{TC} se atiende por \emph{polling} de \texttt{STATUS\_TC} y \textbf{no}
por interrupci\'on; la ISR maneja \'unicamente muestras; y \texttt{CLEAR\_FLAGS}
est\'a prohibido durante \texttt{SAMPLING}.}

\subsection{Resultado medido de la auto-calibraci\'on}

\begin{figure}[!ht]
\centering
\includegraphics[width=0.92\textwidth]{autocalibration_stage_outputs_5_to_20s.png}
\caption{Trayectorias de calibraci\'on de las cuatro etapas anal\'ogicas
capturadas con osciloscopio. La correcci\'on visible abarca unos
\SI{7}{\second}. Cada etapa se calibra en orden y su salida queda dentro de la
banda muerta antes de pasar a la siguiente.}
\label{fig:calib-traj}
\end{figure}

Sobre \textbf{dos placas ensambladas independientemente} y \textbf{ocho
configuraciones placa--ganancia}, con $N=10$ reinicios apareados por
configuraci\'on:

\begin{center}\small
\begin{tabularx}{\textwidth}{@{}l X@{}}
\toprule
\emph{Offset} residual a la salida LP & de \SIrange{2.9}{45.8}{\percent} del fondo de escala a \SIrange{0.4}{3.9}{\percent} \\
Observaciones individuales calibradas & entre \SIrange{0.1}{5.2}{\percent} \\
Prueba de signos de una cola, etapa LP & rechaza la hip\'otesis nula en \textbf{las ocho} configuraciones (10 de 10 reinicios, $p\approx10^{-3}$ exacto); \textbf{sobrevive a la correcci\'on de Bonferroni} \\
Ajuste del controlador & \textbf{\'unica configuraci\'on} para todas las placas y ganancias, sin reajuste por placa ni por ganancia \\
Tiempo & correcci\'on visible \textasciitilde{}\SI{7}{\second}; completitud informal \textasciitilde{}\SI{10}{\second}; ventanas l\'imite configuradas \SI{57.6}{\second} \\
\bottomrule
\end{tabularx}
\end{center}

\dato{\textbf{El \'unico vencimiento de tiempo l\'imite observado en banco} fue
el de una placa alimentada mientras su ge\'ofono estaba en movimiento: el
transitorio de alterna viola la hip\'otesis cuasi-est\'atica que la calibraci\'on
\emph{foreground} requiere. Es una limitaci\'on entendida, no una falla.}

\subsection{Balance de la etapa}

\begin{center}\small
\begin{tabularx}{\textwidth}{@{}l X@{}}
\toprule
\textbf{Se qued\'o} & Controlador PI \'unico con \emph{feedforward}, banda muerta proporcional y anti-\emph{windup} direccional; persistencia en EEPROM con CRC-16 indexada por c\'odigo de ganancia; adquisici\'on por DMA con FIR en hardware; promediado por lote; \texttt{superMaquina} en UDB; comandos USB de diagn\'ostico \\
\addlinespace[3pt]
\textbf{Se descart\'o} & Bisecci\'on binaria; servo lento; t\'ermino derivativo del PID; \texttt{Reg\_Select} como gobierno de la captura; \emph{tick} de sistema de \SI{10}{\milli\second} \\
\addlinespace[3pt]
\textbf{Qued\'o abierto} & Deriva de largo plazo y apareamiento entre canales del arreglo completo, sin caracterizar. La evidencia obtenida concierne al \emph{punto de operaci\'on} de la cadena, no al desempe\~no s\'ismico extremo a extremo \\
\bottomrule
\end{tabularx}
\end{center}

\section{Etapa 7 --- El \emph{pipeline} MASW en Python y las campa\~nas de campo}
\label{sec:et7}
\periodo{3 --- 9 de julio de 2026 \quad$\cdot$\quad campa\~nas los d\'ias 3, 7 y 8; 11 \emph{commits}}

Etapa doble: se sali\'o a medir en serio y, en paralelo, se construy\'o de cero
la cadena de post-proceso que convierte esos registros en un perfil $V_S(z)$.
Es la etapa en la que el proyecto deja de ser un sistema de adquisici\'on y
pasa a ser un m\'etodo completo.

\subsection{La geometr\'ia real de las campa\~nas}

Este dato no estaba volcado en ning\'un documento y se reconstruy\'o para este
historial a partir de los \texttt{metadata.json} de cada carpeta de captura.

\begin{figure}[!ht]
\centering
\includegraphics[width=\textwidth]{geometria_campanas.png}
\caption{Posiciones de receptor efectivamente ocupadas en cada sitio,
reconstruidas desde los metadatos de campo. El \'area del marcador es
proporcional a la cantidad de capturas en esa posici\'on.}
\label{fig:geometria}
\end{figure}

\seleccionado{\textbf{El arreglo multicanal no existe f\'isicamente: se
sintetiza.} Cada captura tiene un \'unico ge\'ofono m\'as el canal de martillo
(\texttt{n\_slaves}~$=2$, roles \texttt{hammer}\,$+$\,\texttt{geo}). El
\emph{gather} multicanal se construye repitiendo el disparo en el mismo origen
y \emph{moviendo el receptor} de posici\'on en posici\'on. Es la forma correcta
de hacer MASW cuando se tienen dos nodos y no veinticuatro, y explica por qu\'e
la sincronizaci\'on y la repetibilidad del disparo fueron tan centrales: sin
repetibilidad entre disparos, el arreglo sintetizado directamente no cierra.}

\begin{center}\small
\begin{tabularx}{\textwidth}{@{}l r r r X@{}}
\toprule
\textbf{Sitio} & \textbf{Pos.} & \textbf{Rango} & \textbf{Capturas} & \textbf{$F_s$ presentes} \\
\midrule
Canchiga    & 18 & \SIrange{10}{50}{\meter} & 74 & \SI{1020}{\hertz} y \SI{2929}{\hertz} \\
Canchita    & 20 & \SIrange{10}{50}{\meter} & 87 & \SI{1020}{\hertz} y \SI{2929}{\hertz} \\
Canchita 2  & 10 & \SIrange{10}{30}{\meter} & 13 & \SI{1020}{\hertz} \\
\bottomrule
\end{tabularx}
\end{center}

\dato{\textbf{Par\'ametros de la geometr\'ia:} fuente (martillo) en el origen,
espaciamiento $\Delta x = \SI{2}{\meter}$, \emph{offset} m\'inimo
\SI{10}{\meter}, apertura del arreglo $L = \SI{40}{\meter}$
(\SI{20}{\meter} en Canchita~2).}

Contrastando contra los criterios de dise\~no que se hab\'ian aprendido en la
Etapa~1:

\begin{center}\small
\begin{tabularx}{\textwidth}{@{}l l X@{}}
\toprule
\textbf{Criterio} & \textbf{Valor} & \textbf{Lectura} \\
\midrule
\emph{Aliasing} espacial & $\lambda_{\min} = 2\Delta x = \SI{4}{\meter}$ &
\textbf{Cumplido abajo, violado arriba.} La longitud de onda m\'as corta en los
\emph{picks} es de \SI{2}{\meter}, la mitad del l\'imite: todo \emph{pick} con
$\lambda < \SI{4}{\meter}$ est\'a \emph{aliased} \\
\addlinespace[2pt]
\emph{Offset} m\'inimo & \SI{10}{\meter} &
Razonable frente al efecto de campo cercano para las frecuencias de trabajo \\
\addlinespace[2pt]
Apertura & $L = \SI{40}{\meter}$ &
\textbf{Es la restricci\'on activa.} La resoluci\'on en baja frecuencia est\'a
acotada por la apertura, y los \emph{picks} llegan a
$\lambda = \SI{130}{\meter}$, m\'as de tres veces $L$ \\
\bottomrule
\end{tabularx}
\end{center}

\pendiente{\textbf{Los \emph{picks} de alta frecuencia hay que revisarlos.} Con
$\Delta x = \SI{2}{\meter}$, cualquier longitud de onda menor que
\SI{4}{\meter} viola el muestreo espacial. Con la velocidad de fase que se
observa en el extremo alto ($\approx\SI{60}{\meter\per\second}$), eso ocurre
por encima de $f \approx c/\lambda_{\min} \approx \SI{15}{\hertz}$ ---y es
justamente la banda donde la curva se aplana en
\SI{60}{\meter\per\second}. Un modo r\'apido con \emph{aliasing} aparece como
un modo lento: la meseta puede ser un artefacto de la geometr\'ia y no una capa
blanda real. Es la \'unica manera de reconciliar la curva con el criterio del
Cap.~3 de Foti, y afecta a m\'as de la mitad de los 51 \emph{picks}.}

\pendiente{\textbf{Esta es la otra observaci\'on cr\'itica.} Con
$L = \SI{40}{\meter}$, la profundidad con soporte geom\'etrico real es del orden
de $L/2 = \SI{20}{\meter}$; el modelo invertido, en cambio, se extiende hasta
\SI{82.6}{\meter} de base. Las capas por debajo de \textasciitilde{}\SI{20}{\meter}
\emph{no est\'an sostenidas por la apertura del arreglo} y deben leerse como
extrapolaci\'on del modelo, no como medici\'on. Para el objetivo declarado de
\SI{50}{\meter} hace falta un arreglo m\'as largo ---y esto es una conclusi\'on
accionable: la extensi\'on de banda hacia abajo ya est\'a resuelta por
electr\'onica, pero la \emph{apertura} sigue siendo el cuello de botella.}

\subsection{El dato crudo, tal cual sali\'o del arreglo}

\begin{figure}[!ht]
\centering
\includegraphics[width=\textwidth]{gather_sintetizado.png}
\caption{\emph{Gather} ensamblado directamente desde los binarios crudos
(\texttt{raw\_f32le.bin}) de la campa\~na Canchita, una traza por posici\'on,
alineadas por el pico del canal de martillo de cada disparo. Cada traza est\'a
normalizada a su propio m\'aximo. Derecha: relaci\'on se\~nal-ruido aparente,
calculada como RMS de la ventana posterior al impacto contra RMS de la ventana
previa del mismo registro. \emph{Verificaci\'on cruzada}: la llegada en la traza
de \SI{10}{\meter} ronda los \SI{200}{\milli\second}, o sea
$\approx\SI{50}{\meter\per\second}$ aparente, consistente con la velocidad
Rayleigh esperable del perfil invertido
($0{,}92\cdot V_{S,\text{sup}} = 0{,}92\cdot\SI{62.9}{} \approx
\SI{58}{\meter\per\second}$). El dato crudo y la inversi\'on concuerdan por
caminos independientes.}
\label{fig:gather}
\end{figure}

Dos cosas se leen directamente de la Fig.~\ref{fig:gather} y no estaban
registradas en ning\'un lado:

\dato{\textbf{El tren dispersivo se ve a simple vista} entre \SI{10}{} y
\SI{26}{\meter}: paquete de baja frecuencia con \emph{moveout} normal, que es
exactamente lo que el m\'etodo necesita.}

\dato{\textbf{La mitad lejana del arreglo est\'a limitada por ruido.} SNR
aparente mediana de \SI{18.9}{\decibel} entre \SI{10}{} y \SI{26}{\meter}
(rango \SIrange{12.7}{25.8}{\decibel}) contra \SI{4.9}{\decibel} entre
\SI{30}{} y \SI{50}{\meter} (rango \SIrange{-2.1}{12.0}{\decibel}): hay
posiciones lejanas \emph{por debajo de \SI{0}{\decibel}}. Se nota adem\'as un
salto de car\'acter entre \SI{26}{} y \SI{30}{\meter}, consistente con que esas
posiciones provienen de jornadas distintas.}

\seleccionado{Esto refuerza la conclusi\'on sobre la apertura desde otro
\'angulo: no alcanza con alargar el arreglo, hay que alargarlo \emph{con
energ\'ia suficiente}. O sea, el l\'imite pr\'actico para llegar a
\SI{50}{\meter} es tanto geom\'etrico como de la fuente de impacto ---lo que
justifica directamente el trabajo del martillo de leva (Etapa~10).}

\subsection{Lo que se implement\'o en Python}

El 7 de julio el repositorio de interfaces sum\'o \textasciitilde{}4600 l\'ineas
de una sola vez:

\begin{center}\small
\begin{tabularx}{\textwidth}{@{}l X@{}}
\toprule
\textbf{M\'odulo} & \textbf{Qu\'e hace} \\
\midrule
\texttt{masw\_dispersion.py} & Del \emph{gather} sintetizado a la imagen de dispersi\'on $(f,c)$ por \emph{phase-shift}. Es la reimplementaci\'on propia del \texttt{p03} de MATLAB de la Etapa~1 \\
\texttt{masw\_inversion.py} & Inversi\'on de la curva a perfil $V_S(z)$ mediante \texttt{evodcinv} sobre \texttt{disba} ---algoritmo evolutivo, elegido por la no unicidad del problema (Cap.~6 de Foti) \\
\texttt{review\_field\_data.py} & Aplicaci\'on de revisi\'on: inspeccionar, filtrar, alinear y aceptar o rechazar registros antes de procesarlos \\
\bottomrule
\end{tabularx}
\end{center}

\subsubsection{Las doce fases de pulido del revisor de campo}

El 8 de julio el revisor recibi\'o doce fases sucesivas. Lo que qued\'o
funcionalmente:

\begin{itemize}[itemsep=0.2em]
\item ordenamiento de trazas por amplitud pico a pico;
\item banco de filtros y pesta\~na de \'enfasis;
\item \textbf{alineaci\'on entre tandas}, necesaria porque las capturas de una
misma posici\'on provienen de disparos distintos;
\item \emph{waterfall} y an\'alisis $f$--$k$;
\item selecci\'on (\emph{picking}) de la curva de dispersi\'on sobre la imagen;
\item \textbf{mezcla de tasas de muestreo}: los registros a \SI{2929}{\hertz} y
a \SI{1020}{\hertz} se combinan por \texttt{resample\_poly}, porque las
campa\~nas del 3 y del 7--8 de julio se tomaron con configuraciones distintas;
\item \emph{deduplicaci\'on} por captura ---el mismo disparo pod\'ia quedar
guardado dos veces;
\item \textbf{polaridad fija}: ge\'ofono no invertido, martillo invertido.
\end{itemize}

\subsubsection{Grupos de dispersi\'on y multimodalidad (9 de julio)}

Se implementaron \textbf{grupos de dispersi\'on independientes}: se definen $N$
regiones-pol\'igono sobre la imagen y cada una produce su propia curva, con
combinaci\'on ponderada de las im\'agenes. Esto es lo que habilita el tratamiento
\textbf{multimodal}: el modo fundamental y los modos superiores se separan como
grupos distintos, se editan a mano si hace falta, y se invierten juntos.

\seleccionado{Invertir con modos superiores adem\'as del fundamental es la
manera est\'andar de atacar la no unicidad: cada modo pesa distinto las capas,
as\'i que usarlos en conjunto restringe el espacio de soluciones mucho m\'as que
el fundamental solo.}

\subsection{En paralelo, en el firmware}

\begin{center}\small
\begin{tabularx}{\textwidth}{@{}l X@{}}
\toprule
\textbf{Cambio} & \textbf{Motivo} \\
\midrule
$F_s$ de \SI{2929}{\hertz} a \SI{1020}{\hertz} (7 jul) & Ajuste de compromiso entre resoluci\'on temporal y volumen de datos por captura \\
\addlinespace[2pt]
Rangos de ADC seleccionables desde la web: primero 2 ($\pm\SI{2.5}{\volt}$ / $\pm\SI{0.512}{\volt}$), luego 4 (8 jul) & Permite adaptar el rango a la distancia fuente--receptor sin recompilar. Todas validadas en hardware con el comando de banco \texttt{range N} \\
\addlinespace[2pt]
Correcci\'on del bug de escala del ADC & Los conteos por volt no correspond\'ian al rango activo \\
\addlinespace[2pt]
\texttt{Timer\_3} dedicado como \emph{watchdog} de captura & Detecta capturas que no progresan \\
\addlinespace[2pt]
\textbf{\emph{Circuit-breaker} anti-tormenta de IRQ} & Se observ\'o en placa, con pyOCD, una tormenta de interrupciones que dejaba al nodo inservible. El breaker la detecta y se auto-recupera \\
\addlinespace[2pt]
Techo de captura del esclavo de 512 a 360 lotes & Regresi\'on encontrada en campo \\
\addlinespace[2pt]
Indicador de RSSI en la web (8 jul) & Diagnosticar en terreno si un nodo est\'a en el borde del alcance, antes de perder una tanda \\
\bottomrule
\end{tabularx}
\end{center}

\subsection{Resultado: el perfil $V_S(z)$}

\begin{figure}[!ht]
\centering
\includegraphics[width=0.92\textwidth]{waterfall_promedios.png}
\caption{\emph{Waterfall} de los promedios de Canchita, tal como lo produce el
revisor de campo tras aceptar los registros: 21 promedios sobre 598 muestras
revisadas. Es la entrada real del c\'alculo de dispersi\'on ---la
Fig.~\ref{fig:gather} muestra el mismo arreglo sin promediar ni filtrar.}
\label{fig:waterfall}
\end{figure}

\begin{figure}[!ht]
\centering
\includegraphics[width=\textwidth]{masw_resultado.png}
\caption{Resultado de la campa\~na Canchita. Izquierda: 51 \emph{picks} de la
curva de dispersi\'on contra la curva te\'orica del modelo invertido
(desajuste declarado \SI{2.55}{\percent}; RMS observado-te\'orico
\SI{22}{\meter\per\second}). Derecha: perfil $V_S(z)$ de nueve capas m\'as
semiespacio. La l\'inea de trazos marca $L/2$, hasta donde la apertura del
arreglo da soporte geom\'etrico.}
\label{fig:masw}
\end{figure}

\dato{\textbf{N\'umeros del resultado (Canchita):}
banda \SIrange{5.11}{29.75}{\hertz};
velocidad de fase observada \SIrange{60}{668}{\meter\per\second};
longitudes de onda \SIrange{2}{130.6}{\meter};
perfil de \SI{62.9}{\meter\per\second} en superficie a
\SI{1011.5}{\meter\per\second} en el semiespacio;
desajuste \SI{2.55}{\percent}.
Del total de 598 muestras revisadas se aceptaron 21 promedios y se
descartaron 9 capturas por criterio expl\'icito.}

\pendiente{El perfil no tiene todav\'ia contraste externo (SPT o ensayo
invasivo). El camino m\'as accesible es la correlaci\'on $V_S$--SPT reportada
para Per\'u y otros contextos regionales, ya en la bibliograf\'ia. Sin ese
contraste, el resultado demuestra que \emph{la cadena cierra de extremo a
extremo}, que no es poco, pero no que el perfil sea correcto.}

\subsection{Balance de la etapa}

\begin{center}\small
\begin{tabularx}{\textwidth}{@{}l X@{}}
\toprule
\textbf{Se qued\'o} & \emph{Pipeline} propio de dispersi\'on e inversi\'on; inversi\'on evolutiva con \texttt{evodcinv}/\texttt{disba}; grupos de dispersi\'on multimodo; revisor de campo con aceptaci\'on expl\'icita de registros; mezcla de $F_s$ por \texttt{resample\_poly}; rangos de ADC conmutables; \emph{watchdog} y \emph{circuit-breaker} de captura; RSSI en la interfaz \\
\addlinespace[3pt]
\textbf{Se descart\'o} & $F_s$ \'unica fija; rango de ADC \'unico; techo de 512 lotes por captura \\
\addlinespace[3pt]
\textbf{Qued\'o abierto} & Apertura del arreglo insuficiente para el objetivo de \SI{50}{\meter}; falta validaci\'on externa del perfil \\
\bottomrule
\end{tabularx}
\end{center}

\section{Etapa 8 --- Auditor\'ia pre-campo}
\label{sec:et8}
\periodo{11 --- 12 de julio de 2026 \quad$\cdot$\quad 2 d\'ias, 33 \emph{commits}}

Dos d\'ias enteros dedicados exclusivamente a \emph{tratar de romper el
sistema antes de que lo rompiera el campo}. Es, en retrospectiva, la etapa
mejor invertida del proyecto: encontr\'o tres defectos que habr\'ian arruinado
campa\~nas completas y que no se manifestaban en las pruebas cortas de
laboratorio.

\begin{figure}[ht]
\centering
\includegraphics[width=\textwidth]{flujo_adquisicion_masw.png}
\caption{Flujo que la auditor\'ia busc\'o preservar de extremo a extremo. Se distinguen
integridad de adquisici\'on, reconstrucci\'on y validaci\'on geot\'ecnica porque no son el
mismo criterio de aceptaci\'on.}
\end{figure}

\subsection{Por qu\'e hac\'ia falta}

Todo lo que se hab\'ia validado hasta ese momento eran capturas cortas. Salir a
campo implica sesiones largas, reconexiones, cortes de energ\'ia, operaci\'on por
interfaz web bajo sol y con una sola oportunidad por sitio. Los modos de falla
que importan en esas condiciones son distintos de los que aparecen en el
banco.

\subsection{La jerarqu\'ia de memoria, validada}

Primero hubo que confirmar que el camino de datos completo funcionaba en las
tres escalas previstas:

\begin{center}\small
\begin{tabularx}{\textwidth}{@{}l X l@{}}
\toprule
\textbf{Nivel} & \textbf{D\'onde vive la captura} & \textbf{Alcance} \\
\midrule
1 & RAM del ESP32 & capturas cortas \\
2 & Encadenado hacia el PSoC & capturas medias \\
3 & \textbf{Tarjeta SD, gestionada por el PSoC} (FatFs, \texttt{GEOLAST.BIN}, comandos \texttt{0xBC}--\texttt{0xBF}) & capturas largas \\
\bottomrule
\end{tabularx}
\end{center}

\seleccionado{La SD vive en el \textbf{PSoC}, no en el ESP. Es coherente con la
divisi\'on de responsabilidades del sistema: el PSoC posee la se\~nal y su
temporizado, el ESP transporta. Se agreg\'o el comando de banco
\texttt{sdread N} para leer bancos directamente y poder auditar lo escrito.}

\subsection{La matriz de aceptaci\'on E1--E18}

\begin{center}\small
\begin{tabularx}{\textwidth}{@{}l X l@{}}
\toprule
\textbf{Id} & \textbf{Prueba} & \textbf{Resultado} \\
\midrule
E1 & RAM chica por radio (\texttt{-n 4}) & PASS 120/120 \\
E2 & SD chica por radio (\texttt{-n 600}, \SI{6.9}{\second}) & PASS 18\,000/18\,000 \\
E3 & SD mediana por radio (\texttt{-n 3000}, \SI{34.6}{\second}) & PASS 90\,000/90\,000 \\
E4 & Log paralelo durante E3 & \SI{3.6}{\percent} de lecturas SD sobre \SI{250}{\milli\second} (defecto H3) \\
E5 & \textbf{Aceptaci\'on de 10 minutos} (52\,080 lotes) & \textbf{FAIL --- 0 muestras} \\
E5b/E5c & Repetici\'on con el fix F1 & \textbf{PASS extremo a extremo} \\
E6/E6b & Decimaci\'on 2 y FIR por radio & PASS 3000/3000 \\
E7--E9 & Comandos \texttt{0xBE}, \texttt{0xAD}, USB estricto & PASS (33/33 en USB) \\
E10 & Interfaz en navegador & PASS \\
E12 & \textbf{Reconexi\'on WebSocket cortada durante el volcado} & PASS exacto \\
E13 & STOP en vivo, en RUNNING y en DUMPING & PASS \\
E14 & Operaci\'on completa desde la interfaz & PASS \\
E15 & Captura SD m\'axima, 60\,000 lotes & PASS \\
E16 & Transiciones de rango y decimaci\'on & PASS 49/49 \\
E17 & Ciclo de energ\'ia y recuperaci\'on & PASS 30/30 \\
E18 & \emph{Mini-soak} & PASS 5/5 \\
\bottomrule
\end{tabularx}
\end{center}

\dato{\textbf{Resultado de la prueba de aceptaci\'on final (E5c):}
52\,080/52\,080 lotes escritos en SD y 1\,562\,400/1\,562\,400 muestras
transferidas por radio, sobre una captura continua de 10 minutos. Cero
p\'erdidas. Veredicto registrado: \emph{sistema del banco listo para campo}.}

\subsection{Los tres defectos cr\'iticos}

\subsubsection{F1 --- el \emph{watchdog} destru\'ia las capturas largas}

El defecto que produjo el ``10 minutos $=$ 0 muestras'', y el m\'as instructivo
de todos: \textbf{el mecanismo de protecci\'on era la causa de la falla}.

El \emph{watchdog} \texttt{SYNC\_STALL} exist\'ia para detectar que una captura
se hab\'ia trabado. Su margen era fijo, no proporcional a la duraci\'on de la
captura, as\'i que en sesiones largas ---donde escribir en SD demora
leg\'itimamente--- disparaba solo, y su acci\'on de recuperaci\'on era
\emph{re-disparar} \texttt{START\_NOW}. Cada re-disparo destru\'ia la sesi\'on en
curso y arrancaba otra, que a su vez volv\'ia a disparar el \emph{watchdog}. El
resultado eran diez minutos de trabajo y cero muestras.

\seleccionado{\textbf{Fix:} margen proporcional a la duraci\'on esperada, y
aborto limpio en vez de re-disparo. La lecci\'on transferible es que un
\emph{watchdog} cuya acci\'on correctiva reinicia el trabajo puede convertir una
lentitud tolerable en una p\'erdida total, y que \textbf{esa clase de defecto
solo aparece en pruebas de duraci\'on realista}.}

\subsubsection{F9 --- START romp\'ia el nodo}

\dato{Marcado como cr\'itico: \emph{START romp\'ia el nodo y lo dejaba inservible
hasta reset f\'isico, siempre, con cualquier $N$}. Causa ra\'iz: una condici\'on
de carrera entre el reintento de PRESTART del maestro y un estado ARMED del
PSoC que no se anunciaba. El maestro reintentaba mientras el PSoC ya hab\'ia
transicionado en silencio, y el nodo quedaba en un estado inconsistente del
que no sal\'ia solo.}

Es la contracara del logro de la Etapa~4: el mismo mecanismo de
sincronizaci\'on que resolvi\'o el arranque coordinado ten\'ia un camino de
reintento sin cubrir.

\subsubsection{F5 --- volcado colgado por socket zombi}

La herramienta de captura por WebSocket se colgaba indefinidamente si el
socket quedaba medio abierto: sin datos, sin error, sin fin. Se corrigi\'o con
un \emph{watchdog} de silencio y endurecimiento previo a la fase de datos, y
se valid\'o cortando la conexi\'on a prop\'osito durante el volcado (E12), que
reconect\'o y termin\'o bien.

\subsection{Defectos menores documentados y no corregidos}

Registrarlos vale tanto como corregirlos, porque evita volver a investigarlos:

\begin{center}\small
\begin{tabularx}{\textwidth}{@{}l X l@{}}
\toprule
\textbf{Id} & \textbf{Descripci\'on} & \textbf{Estado} \\
\midrule
H3 & \SI{3.6}{\percent} de lecturas SD superan \SI{250}{\milli\second} (picos de \SIrange{1.2}{1.5}{\second}) & Mitigado: \emph{timeout} \SI{250}{}$\to$\SI{450}{\milli\second}; el reintento siempre salva \\
F2 & Los ACK muestran \texttt{value \& 0xFF} (600 se ve como 88) & Cosm\'etico, documentado \\
W3 & STOP durante el volcado: no hay re-volcado desde la interfaz (los datos quedan en la SD) & Documentado \\
W5 & La vista previa de $N$ usa el campo de entrada y no el ACK & Documentado \\
H4 & \emph{Prefetch} de bloque en el PSoC dar\'ia un volcado \textasciitilde{}$2\times$ m\'as r\'apido & Mejora identificada, post-campo \\
W6/W7 & M\'aximo de RAM en 513 lotes por redondeo hacia arriba; paso del campo de duraci\'on & Corregidos \\
\bottomrule
\end{tabularx}
\end{center}

\subsection{Balance de la etapa}

\begin{center}\small
\begin{tabularx}{\textwidth}{@{}l X@{}}
\toprule
\textbf{Se qued\'o} & Jerarqu\'ia de memoria RAM$\to$encadenado$\to$SD validada en hardware; SD gestionada por el PSoC con FatFs; matriz E1--E18 como procedimiento de aceptaci\'on repetible; suite de \emph{scripts} de regresi\'on (USB, SD m\'axima, transiciones de rango, recuperaci\'on por reset); \emph{watchdog} de silencio en las herramientas \\
\addlinespace[3pt]
\textbf{Se descart\'o} & \emph{Watchdog} con margen fijo; re-disparo de \texttt{START\_NOW} como acci\'on de recuperaci\'on; criterios de prueba que se bloqueaban por contadores benignos (F8) \\
\addlinespace[3pt]
\textbf{Aprendizaje} & Los defectos que importan en campo son de \emph{duraci\'on} y de \emph{reconexi\'on}, y no aparecen en pruebas cortas. Tres de ellos habr\'ian costado campa\~nas enteras \\
\bottomrule
\end{tabularx}
\end{center}

\section{Etapa 9 --- Art\'iculo URUCON y caracterizaci\'on medida de la cadena}
\label{sec:et9}
\periodo{15 --- 23 de julio de 2026 \quad$\cdot$\quad 5 d\'ias de trabajo}

Dos l\'ineas: escribir y someter a revisi\'on el art\'iculo sobre la
auto-calibraci\'on, y ---en paralelo--- \textbf{medir de verdad} la cadena
anal\'ogica con generador y osciloscopio, en vez de confiar en el modelo.

\begin{figure}[ht]
\centering
\includegraphics[width=\textwidth]{identificacion_compensador_medida.png}
\caption{Identificaci\'on experimental hist\'orica del compensador. Los puntos medidos
documentan la respuesta observada; la curva de modelo superpuesta es provisional porque
el primer modelo de operacional confundi\'o GBW con polo dominante.}
\end{figure}

\begin{figure}[ht]
\centering
\includegraphics[width=\textwidth]{geofono_vs_cadena_medida.png}
\caption{Comparaci\'on hist\'orica entre respuesta del ge\'ofono y cadena medida. Se
conserva como diagn\'ostico, no como validaci\'on vigente de la banda compensada hasta
repetir la s\'intesis y la medici\'on.}
\end{figure}

\subsection{El art\'iculo: \emph{Automatic DC-Offset Calibration for a Low-Cost
Geophone Front End}}

\begin{center}\small
\begin{tabularx}{\textwidth}{@{}l X@{}}
\toprule
15 jul & Reescritura de las Secciones III y IV con datos experimentales reales (dos placas, barrido de ganancia completo). Se agrega el \emph{flag} de firmware \texttt{CAL\_PI\_FORCE\_MIN\_DEADBAND} para el experimento decisivo sobre el estancamiento del PGA \\
\addlinespace[2pt]
15 jul & \textbf{Revisi\'on editorial simulada}: se somete el propio art\'iculo a un panel de cinco revisores ficticios antes de mandarlo \\
\addlinespace[2pt]
16 jul & Respuesta a esa revisi\'on: siete referencias nuevas, tabla de resultados ampliada a las cuatro etapas (no solo LP), ecuaciones de banda muerta con justificaci\'on expl\'icita, y \textbf{una correcci\'on de honestidad} sobre qu\'e es realmente la condici\'on ``Default'' \\
\addlinespace[2pt]
23 jul & Reuni\'on con el tutor; compresi\'on de 6 a 5 p\'aginas; figura GEO$\to$ADC lado a lado nueva. Panel de revisi\'on simulado da \emph{Minor Revision}, $\approx$65/100 \\
\bottomrule
\end{tabularx}
\end{center}

\seleccionado{Vale registrar el m\'etodo: \textbf{someter el trabajo a una
revisi\'on adversarial antes de exponerlo}. Lo que produjo no fue cosm\'etico
---la ampliaci\'on de la tabla a las cuatro etapas y la correcci\'on sobre la
condici\'on ``Default'' son cambios de fondo que salieron de esa cr\'itica.}

Del mismo per\'iodo es el cierre de una deuda de nomenclatura: el componente
\texttt{GEO\_ADDER} se renombr\'o a \texttt{GEO\_SUM} en todo el firmware
(esquem\'atico, calibraci\'on, tablas por etapa), con recalibraci\'on de la
ganancia del PI ($\times2{,}63$) y banda muerta configurable por etapa.

\subsection{El protocolo de barridos: c\'omo se midi\'o la cadena}

Este es el documento t\'ecnico m\'as riguroso que produjo el proyecto
(\texttt{PROTOCOLO\_BARRIDOS\_CADENA\_ANALOGICA.md}). Registra c\'omo medir la
respuesta en frecuencia de las cuatro etapas con un generador Agilent 33220A y
un osciloscopio Tektronix TDS2004B:

\begin{itemize}[itemsep=0.2em]
\item CH1 $=$ salida PGA, CH2 $=$ BP, CH3 $=$ sumador (compensador), CH4 $=$ LP;
\item \texttt{SYNC} del generador $\to$ \texttt{EXT TRIG} del osciloscopio, de
modo que el flanco de \emph{sync} marque el comienzo del barrido;
\item \textbf{sondas en $10\times$ obligatoriamente}: en $1\times$ la
capacitancia de la sonda es comparable con el $C_f = \SI{177}{\pico\farad}$ del
nodo BP y lo carga de forma apreciable ---o sea, \emph{el instrumento
modificar\'ia el circuito que quiere medir};
\item nueve bandas de barrido logar\'itmico desde
\SIrange{100}{1}{} m\SI{}{\hertz}--\SI{}{\hertz} hasta
\SIrange{20}{200}{\kilo\hertz}, cada una con su base de tiempo;
\item tres repeticiones por banda cuando sea posible;
\item \textbf{control previo de sondas}: conectar los cuatro canales al mismo
seno y verificar que las razones den $1\angle 0^{\circ}$, para detectar una
sonda descompensada \emph{antes} de atribuirle el error al circuito.
\end{itemize}

\dato{\textbf{Precauci\'on documentada sobre la etapa SUM.} En la banda donde
$BP/PGA \approx -1$, las dos corrientes de entrada del sumador se cancelan en
torno al \SI{95}{\percent}. Eso hace que el cociente medido sea
extremadamente sensible al ruido y a los valores reales de los resistores. La
recomendaci\'on registrada es \emph{abrir temporalmente una de las dos ramas}
para identificar la etapa sin la cancelaci\'on, y usar la medici\'on de cuatro
canales solo como validaci\'on en la configuraci\'on final.}

\subsection{Tratamiento de datos: lo que se hizo bien}

El script de an\'alisis (\texttt{analizar\_barridos\_cadena\_analogica.m},
\texttt{AnalisisCircuito/}) incorpora varias decisiones que conviene no perder:

\begin{itemize}[itemsep=0.2em]
\item \textbf{Modelo del operacional, no ideal}: se registraron
$A_0=\SI{90}{\decibel}$, GBW $=\SI{8}{\mega\hertz}$,
$R_{in}=\SI{35}{\mega\ohm}$,
$R_{out}=\SI{20}{\ohm}$, $C_{in}=\SI{18}{\pico\farad}$. El \emph{slew rate}
(\SI{4.3}{\volt\per\micro\second}) se trata como l\'imite \emph{no lineal} de
amplitud y no como un polo, que es lo correcto. Sin embargo, otros scripts
introdujeron el GBW de \SI{8}{\mega\hertz} como si fuera el polo dominante; esa
parte del resultado queda invalidada. La relaci\'on $f_p=f_t/A_0$ ya se corrigi\'o
en el c\'odigo, pero falta repetir la s\'intesis y el contraste.
\item \textbf{Cancelaci\'on de red por m\'inimos cuadrados}, la misma t\'ecnica
que en el ESP: se estima la fundamental de \SI{50}{\hertz} y se restan
\emph{solo} los arm\'onicos dentro del pasabanda. Se excluye del ajuste un
entorno de tres ciclos alrededor de cada cruce del barrido \textbf{para no
borrar el est\'imulo}, y se omite la regresi\'on si el registro no contiene al
menos dos ciclos de red.
\item \textbf{Recuperaci\'on parcial de capturas saturadas}: una saturaci\'on ya
no descarta la captura entera. Se excluye un ciclo alrededor de cada recorte y
los tramos lineales continuos restantes entran como submuestras
independientes. En la tanda calibrada esto recuper\'o 22 combinaciones
captura-etapa de 9 capturas saturadas.
\item \textbf{Deduplicaci\'on por SHA-256}: se detect\'o que \texttt{XD\_ALL0006}
era copia exacta de \texttt{XD\_ALL0005} en los cuatro canales y se elimin\'o.
\item \textbf{Orden ampliado con criterio}: dos polos extra por cada
operacional en el camino, aumentando tambi\'en el grado del numerador para
conservar el grado relativo (BP 4/3, compensador 5/4, LP 4/2, PGA$\to$ADC
11/8).
\end{itemize}

\subsection{Un conjunto de datos completo, descartado}

\descartado{\textbf{\texttt{Osciloscopio\_descartado\_2026-07-20\_BP\_43k}} ---
una tanda completa de barridos, con sus 156 archivos y m\'as de 20 gr\'aficos ya
producidos, qued\'o archivada como descartada. Motivo: la rama BP ten\'ia
instalado \SI{8200}{\ohm} donde el dise\~no requiere \SI{7501.9}{\ohm}, de modo
que la etapa no estaba realizando la compensaci\'on buscada. Se conserva en el
repositorio en vez de borrarse, con el motivo en el nombre de la carpeta.}

\seleccionado{Conservar el conjunto descartado \emph{con el motivo en el
nombre} es la decisi\'on correcta: permite volver a \'el si aparece una duda, y
evita que alguien lo reprocese sin saber por qu\'e no vale.}

\subsection{La tanda calibrada del 21 de julio}

Se rehicieron los barridos con la rama BP corregida a
$\SI{6.8}{\kilo\ohm} + $ potenci\'ometro. Resultados de identificaci\'on:

\begin{center}\small
\begin{tabularx}{\textwidth}{@{}l r r r X@{}}
\toprule
\textbf{Etapa} & \textbf{Pts.} & \textbf{Coh.\ med.} & \textbf{RMS ajuste} & \textbf{Estado} \\
\midrule
BP (CH2/CH1) & 85 & 0{,}968 & \SI{6.49}{\decibel} & Parcial: polo inferior fuera de banda, alta frecuencia no ideal \\
COMP (CH3/CH1) & 71 & 0{,}964 & \SI{0.38}{\decibel} & Identificaci\'on buena; el pot no es estimable de forma confiable \\
LP (CH4/CH3) & 33 & 0{,}945 & \SI{0.60}{\decibel} & Confiable en la banda \'util \\
PGA$\to$ADC (CH4/CH1) & 61 & \textbf{0{,}995} & \SI{0.20}{\decibel} & Confiable; din\'amica subs\'onica no observable \\
\bottomrule
\end{tabularx}
\end{center}

\dato{Se extendi\'o la medici\'on hasta \SI{10}{\milli\hertz} con la captura
\texttt{XD\_ALL0017} (\SIrange{0.01}{0.05}{\hertz}), libre de saturaci\'on, que
aporta 14 puntos por transferencia con coherencia mediana entre
0{,}988 y 0{,}9998.}

\subsubsection{El hallazgo principal: la muesca no llega a profundidad}

\dato{En $f_0 = \SI{10.2046}{\hertz}$, el modelo completo predice
$CH3/CH1 = \SI{-57.58}{\decibel}$. La medici\'on da $\approx\SI{-31.26}{\decibel}$:
la muesca \textbf{es reproducible} ---dos capturas independientes difieren solo
\SI{0.42}{\decibel}, y dos barridos nuevos concuerdan--- pero queda
\textbf{\SI{26.33}{\decibel} menos profunda} que el objetivo.}

La causa est\'a identificada y es geom\'etrica, no de firmware:

\begin{center}\small
\begin{tabularx}{\textwidth}{@{}l r X@{}}
\toprule
Cursor para cancelaci\'on exacta & \SI{632.556}{\ohm} & \multirow{2}{*}{} \\
Cursor objetivo para $\zeta_0 = 0{,}25$ & \SI{634.539}{\ohm} & \\
\textbf{Separaci\'on entre ambos ajustes} & \textbf{\SI{1.98}{\ohm}} & sobre un potenci\'ometro de recorrido \SI{2}{\kilo\ohm} \\
Posici\'on efectiva estimada por profundidad & $\approx\SI{674}{\ohm}$ & unos \SI{39}{\ohm} por encima del objetivo \\
\bottomrule
\end{tabularx}
\end{center}

\seleccionado{\textbf{El ajuste requerido es de \SI{1.98}{\ohm} sobre un
recorrido de \SI{2000}{\ohm}: una parte en mil.} Esa es la raz\'on real por la
que la compensaci\'on no alcanza su profundidad de dise\~no, y es una conclusi\'on
de dise\~no, no un problema de calibraci\'on: \emph{la sensibilidad del ajuste
est\'a mal escalada}. La conclusi\'on pr\'actica que qued\'o escrita es no seguir
con m\'as barridos anchos, sino recalibrar alrededor de \SI{10.2046}{\hertz} con
mayor sensibilidad vertical en CH3 y repetir un barrido estrecho.}

\dato{\textbf{Lo que s\'i funciona:} el LP se comporta como se espera
(\SI{13.19}{}/\SI{13.20}{\decibel} medidos contra \SI{13.98}{\decibel} del
modelo; \SI{0.435}{\decibel} RMS en toda la banda de ajuste
\SIrange{10.6}{984}{\hertz}), y la cadena PGA$\to$ADC se identifica con
\SI{0.20}{\decibel} y \SI{1.16}{\degree} RMS. \textbf{El circuito es
repetible}; lo que no est\'a logrado es la calibraci\'on fina de la cancelaci\'on.}

\subsection{Lectura resuelta de los tres valores de \texorpdfstring{$\zeta$}{zeta}}

Los archivos de resultados distinguen tres cantidades que se hab\'ian venido
mezclando:
\begin{itemize}[itemsep=0.15em,topsep=0.2em]
\item \texttt{zeta0\_objetivo} $= 0{,}25$, descrito como \emph{``numerador de
Hcompensate''}. Es el amortiguamiento del ge\'ofono, que el numerador del
compensador debe reproducir para cancelar sus polos.
\item \texttt{zeta1\_denominador\_BP}: $980{,}03$ en la tanda descartada del
20/07 y $937{,}40$ en la calibrada del 21/07. Es el amortiguamiento
\emph{nuevo}, el que efectivamente extiende la banda.
\item \texttt{zeta\_realizada\_componentes}: \textbf{$83{,}661$} en la tanda
descartada (con la resistencia de \SI{8200}{\ohm} equivocada) y
\textbf{$0{,}250000000000114$} en la calibrada.
\end{itemize}
Por tanto, \textbf{$83{,}661$ no es un objetivo de dise\~no}: es el valor que
el numerador realizaba por error con la resistencia equivocada. La tanda
calibrada lo lleva a $0{,}25$, mientras que $937{,}40$ permanece como
amortiguamiento del denominador BP. Esta lectura queda adoptada como
autoritativa y corrige la interpretaci\'on anterior del borrador de entrega.

\subsection{Balance de la etapa}

\begin{center}\small
\begin{tabularx}{\textwidth}{@{}l X@{}}
\toprule
\textbf{Se qued\'o} & Protocolo de barridos reproducible; identificaci\'on por \texttt{tfest} con orden justificado; modelo de operacional real; cancelaci\'on de red por m\'inimos cuadrados tambi\'en en el an\'alisis; recuperaci\'on parcial de capturas saturadas; deduplicaci\'on por SHA-256; revisi\'on adversarial simulada como paso previo a publicar \\
\addlinespace[3pt]
\textbf{Se descart\'o} & Tanda completa del 20/07 con BP de \SI{43}{\kilo\ohm} y rama de \SI{8200}{\ohm}; nomenclatura \texttt{GEO\_ADDER}; sondas en $1\times$ \\
\addlinespace[3pt]
\textbf{Qued\'o abierto} & La muesca de compensaci\'on est\'a \SI{26.33}{\decibel} por encima de su profundidad de dise\~no; el ajuste de \SI{1.98}{\ohm} sobre \SI{2}{\kilo\ohm} es impracticable como est\'a y sugiere rehacer el divisor con mayor resoluci\'on \\
\bottomrule
\end{tabularx}
\end{center}

\section{Etapa 10 --- Ruptura arquitect\'onica y portaci\'on a la web}
\label{sec:et10}
\periodo{21 de julio --- 1\textsuperscript{o} de agosto de 2026 \quad$\cdot$\quad 7 d\'ias, \textasciitilde{}40 \emph{commits}}

Etapa de infraestructura. No agrega capacidad de medici\'on, pero cambia c\'omo
se sostiene todo lo dem\'as ---y contiene un pivote de herramienta importante.

\begin{figure}[ht]
\centering
\includegraphics[width=\textwidth]{flujo_adquisicion_masw.png}
\caption{Flujo funcional que gui\'o la separaci\'on entre firmware, datos y aplicaci\'on
web dentro del superproyecto.}
\end{figure}

\subsection{Del monolito al superproyecto (21 de julio)}

El repositorio \'unico se parti\'o. El historial monol\'itico se preserv\'o en
\texttt{Tesis-legacy} y el repositorio \texttt{Tesis} se reinici\'o con un
\textbf{\emph{commit} ra\'iz sin padres} que fija los subm\'odulos como
repositorios independientes.

\begin{center}\small
\begin{tabularx}{\textwidth}{@{}l X@{}}
\toprule
\textbf{Motivo} & \textbf{Consecuencia} \\
\midrule
El firmware, el modelado y los datos ten\'ian ciclos de vida distintos & Cada uno versiona a su ritmo \\
\addlinespace[2pt]
Los datos crudos (\textasciitilde{}29~GiB) envenenaban el historial & LFS con almacenamiento por carpeta y deduplicaci\'on SHA-256, con manifiesto de ruta/tama\~no/\emph{hash} por archivo \\
\addlinespace[2pt]
Un clon completo era impracticable & Se puede clonar solo lo que se necesita \\
\bottomrule
\end{tabularx}
\end{center}

\subsection{Reorganizaci\'on por prop\'osito (24 de julio)}

Tres d\'ias despu\'es, segunda pasada: \texttt{src/} deja de agrupar por lenguaje
y pasa a agrupar \textbf{por prop\'osito}
---\texttt{firmware/}, \texttt{interfaces/}, \texttt{calculos\_modelados/}---
mediante \texttt{subtree split}, con dos repositorios nuevos y dos renombrados.
\texttt{investigacion/} baj\'o de la ra\'iz a \texttt{docs/investigacion}.

\seleccionado{Agrupar por prop\'osito y no por lenguaje es lo correcto ac\'a
porque \emph{el mismo prop\'osito se implementa en varios lenguajes}: hay
interfaces en Python y en MATLAB, y modelado en Python y en MATLAB. Agrupar
por lenguaje separaba cosas que se usan juntas.}

\pendiente{Dos trampas conocidas de esta arquitectura, que ya costaron trabajo:
\textbf{(1)} un subm\'odulo en \emph{detached HEAD} acepta un
\texttt{push origin main} \emph{sin subir nada}, informando \'exito ---hay que
auditar el \emph{gitlink} contra \texttt{ls-remote}; \textbf{(2)}
\texttt{git submodule deinit -f} vac\'ia el \emph{worktree} y destruye lo no
confirmado. Para mover un subm\'odulo hay que editar solo el \texttt{path} en
\texttt{.gitmodules}.}

\subsection{El pivote de interfaz: de PyQt a web}

Hasta julio la revisi\'on de campo viv\'ia en \texttt{field\_review\_app}, una
aplicaci\'on de escritorio PyQt. Entre el 25 de julio y el 29 se port\'o
entera a la web.

\begin{center}\small
\begin{tabularx}{\textwidth}{@{}l X@{}}
\toprule
25 jul & Arranca el \texttt{PORT\_PLAN}: el servidor Python migra de un \texttt{http.server} artesanal a \textbf{FastAPI} con est\'aticos; pesta\~nas de navegaci\'on y alternancia de tema; visor de se\~nal con \textbf{decimado min/max}. En el camino se corrige un bug del propio \emph{gate} automatizado, que daba falsos vencimientos porque un \emph{pipe} de subproceso no ten\'ia lector \\
\addlinespace[2pt]
26 jul & \textbf{Se congela un \emph{snapshot} de referencia de la app PyQt completa} en \texttt{docs/legacy/pyqt\_field\_review/}, con mapa de l\'ineas de la ventana ``Capturas'', \emph{antes} de empezar a portarla \\
\addlinespace[2pt]
29 jul & Portaci\'on masiva hasta alcanzar paridad funcional \\
\bottomrule
\end{tabularx}
\end{center}

\seleccionado{\textbf{Congelar el original antes de portar} es la decisi\'on que
hizo viable el porte: da un objetivo de paridad verificable en vez de una
reescritura ``a ojo''. El decimado min/max, por su parte, es la t\'ecnica
correcta para dibujar millones de muestras en un \emph{canvas} sin perder los
picos ---que en una traza s\'ismica son exactamente lo que importa.}

\subsubsection{Lo que trajo la portaci\'on (29 de julio)}

\begin{itemize}[itemsep=0.2em]
\item \textbf{Cuarentena reversible} (\texttt{deletion.py}) en vez de borrado
f\'isico: un registro rechazado se puede recuperar. Dado que un descarte
err\'oneo de datos de campo es irreparable, es la elecci\'on correcta.
\item \textbf{Cola \texttt{AnalysisJobs}} persistente para los an\'alisis
pesados, que ya no bloquean la interfaz.
\item Exportaciones reproducibles.
\item \textbf{\texttt{smoke\_test.py} con 58 verificaciones} como condici\'on de
avance: el porte no avanzaba si el \emph{gate} no pasaba.
\end{itemize}

\subsection{Otras l\'ineas del per\'iodo}

\textbf{Conectividad de campo.} Arranca el plan del modo ENLACE del maestro y
un \emph{sketch} descartable (\texttt{MasterFoo}) que valida el uso del
\emph{hotspot} del celular como red de campo.

\descartado{\texttt{MasterFoo} era deliberadamente descartable: se escribi\'o
para responder una pregunta (¿funciona el \emph{hotspot}?), se la respondi\'o y
no se lo integr\'o. Registrarlo evita que alguien lo confunda con c\'odigo de
producci\'on.}

\textbf{Martillo de leva} (1\textsuperscript{o} de agosto). Arranca el dise\~no
de una fuente de impacto repetible ---clave para el arreglo sintetizado de la
Etapa~7, donde la repetibilidad entre disparos \emph{es} la
calidad del dato. Se model\'o en dos herramientas independientes: Simscape
Multibody (MATLAB) y un modelo 2D en Python.

\dato{\textbf{Los dos modelos detectaron errores propios distintos}, lo cual es
justamente el argumento para hacerlos por duplicado: el de Simscape encontr\'o
una \emph{interferencia axial entre la leva y el brazo}, y el de Python una
\emph{violaci\'on de conservaci\'on de energ\'ia} en el criterio de suelta.}

\textbf{PCBs.} Entra el proyecto de placa (KiCad, y luego JitX) como
repositorio propio.

\subsection{Balance de la etapa}

\begin{center}\small
\begin{tabularx}{\textwidth}{@{}l X@{}}
\toprule
\textbf{Se qued\'o} & Superproyecto con subm\'odulos por prop\'osito; LFS con dedup SHA-256 y manifiesto; servidor FastAPI; interfaz de revisi\'on en web; decimado min/max; cuarentena reversible; cola \texttt{AnalysisJobs}; \emph{gate} de 58 verificaciones; martillo de leva modelado por duplicado \\
\addlinespace[3pt]
\textbf{Se descart\'o} & Repositorio monol\'itico; organizaci\'on de \texttt{src/} por lenguaje; \texttt{http.server} artesanal; app de escritorio PyQt (congelada como referencia, no borrada); \emph{sketch} \texttt{MasterFoo}; borrado f\'isico de registros \\
\bottomrule
\end{tabularx}
\end{center}

\section{Etapa 11 --- La placa nueva (rama \texttt{cambios-hardware})}
\label{sec:et11}
\periodo{2 --- 12 de agosto de 2026 \quad$\cdot$\quad etapa en curso}

\dato{La bit\'acora diaria termina el 1\textsuperscript{o} de agosto. Esta
secci\'on est\'a reconstruida directamente de los \emph{commits} y de la
documentaci\'on de construcci\'on del firmware, no de la bit\'acora.}

Se dise\~n\'o una placa propia ---hasta ac\'a todo hab\'ia corrido sobre
\emph{kits} de desarrollo y una \emph{carrier} intermedia
(CY8CKIT-059 $+$ ESP32 DevKitC)--- y eso oblig\'o a rehacer el TopDesign del
PSoC y buena parte del enlace entre procesadores. Todo el trabajo vive en la
rama \texttt{cambios-hardware}, deliberadamente separada de \texttt{main}.

\subsection{Cambios en el PSoC}

\begin{center}\small
\begin{tabularx}{\textwidth}{@{}l X@{}}
\toprule
\textbf{Cambio} & \textbf{Motivo y consecuencia} \\
\midrule
\textbf{IDAC en lugar de VDAC} en las cuatro referencias de calibraci\'on &
Se pasa a fijar las referencias por \emph{corriente} en vez de por tensi\'on,
con las globales de placa $R = \SI{30}{\kilo\ohm}$,
$V_{\mathrm{ref}} = \SI{2.048}{\volt}$,
$I_{fs} = \SI{31.875}{\micro\ampere}$ y helpers de conversi\'on
c\'odigo$\leftrightarrow$tensi\'on. \textbf{El LSB queda en
\SI{3.75}{\milli\volt}}, que es la nueva resoluci\'on del actuador del lazo PI
de la Etapa~6 \\
\addlinespace[3pt]
\textbf{La UART queda solo como RX} &
El TX desaparece del TopDesign nuevo \\
\addlinespace[3pt]
\textbf{Subida PSoC$\to$ESP por I\textsuperscript{2}C} (\texttt{psoc\_link\_*}, PSoC maestro) &
\emph{La UART a 115\,200 era el cuello de botella del \emph{stream} de
muestras.} El protocolo de \emph{frames} no cambia: cambia el medio \\
\addlinespace[3pt]
\textbf{\texttt{PGAout}}: nueva etapa de salida del \emph{pipeline} GEO &
Configurable por el comando \texttt{0xA8}, igual que \texttt{PGAgain} y \texttt{PGAvdac} \\
\bottomrule
\end{tabularx}
\end{center}

\seleccionado{\textbf{El principio de la Etapa~4 sobrevive intacto
al cambio de placa}: el medio del enlace cambi\'o por tercera vez
(SPI $\to$ UART $\to$ I\textsuperscript{2}C) sin tocar el protocolo. Que el
formato de \emph{frames} sea independiente del transporte es lo que hizo
barata cada migraci\'on.}

\subsubsection{Trampas del TopDesign nuevo, registradas}

\pendiente{Estas cost\'o encontrarlas y conviene no volver a tropezar:
\begin{itemize}[itemsep=0.15em,topsep=0.2em]
\item El componente se llama \texttt{I2C\_1}, \textbf{no} \texttt{I2C}.
\item \textbf{El TopDesign nuevo no tiene pin LED}: hubo que sacar el
\emph{include} directo de \texttt{LED.h} para que los \emph{guards}
\texttt{CY\_PINS\_LED\_H} hicieran efecto.
\item Los pines de \texttt{SPIp} quedaron \textbf{repartidos en varios
puertos}, as\'i que el \emph{fitter} dej\'o de emitir los macros agregados; para
el CS de la SD se usan los del pin~0.
\end{itemize}
Y sigue vigente la regla del proyecto: \textbf{nunca editar a mano}
\texttt{.cyprj}, \texttt{.cydwr}, \texttt{.cysch} ni \texttt{.cyfit} ---solo
\texttt{.c} y \texttt{.h}. Romper esto ya rompi\'o un \emph{build} una vez.}

\dato{Estado del \emph{build} tras adaptar el firmware al TopDesign nuevo:
\textbf{compila}. Flash \SI{60400}{\byte}, SRAM \SI{51928}{\byte}.}

\subsection{Cambios en el ESP32}

\textbf{El enlace pasa a ser asim\'etrico.} ESP\,$\to$\,PSoC sigue por UART
(porque el PSoC conserva el RX); PSoC\,$\to$\,ESP pasa a I\textsuperscript{2}C
con el PSoC de maestro y el ESP de esclavo en \texttt{PSOC\_I2C\_ADDR}
(GPIO\,21/22). Los bytes entran a un \emph{ring buffer} y \texttt{poll()} los
consume con el mismo \emph{parser} de \emph{frames} de siempre.

\textbf{Sincronizaci\'on: se elimina el cable.} Desaparecen
\texttt{SYNC\_IN\_PIN} y la ISR \texttt{onSyncEdge}. Ya no hay sincronizaci\'on
por cable desde el maestro: \textbf{el arranque llega \'unicamente por ESP-NOW},
y es el propio esclavo el que genera el flanco hacia el PSoC por
\texttt{SYNC\_TO\_PSOC\_PIN}, el \'unico pin de \emph{sync} que queda.

\seleccionado{Esto simplifica el despliegue en campo ---un cable menos por
nodo--- \textbf{sin violar el principio de instante cero por flanco f\'isico}:
el flanco sigue existiendo, solo que ahora lo genera el esclavo localmente al
recibir el disparo por radio. PRESTART/HOT\_WAIT garantiza que todos est\'en
listos; la dispersi\'on temporal entre nodos queda gobernada por la entrega y el
despacho de \texttt{CMD\_START} m\'as el camino GPIO$\to$PSoC.}

\pendiente{Consecuencia directa: \textbf{el presupuesto de \emph{jitter} hay
que rehacerlo} para esta topolog\'ia. Con \emph{sync} por cable, una parte del
error estaba acotada por hardware; ahora todo el error de sincronizaci\'on pasa
por el enlace de radio. Es la verificaci\'on m\'as urgente de esta etapa.}

El firmware incluye un ensayo reproducible para cerrarla:
\texttt{SCOPE\_MULTI} ejecuta 128 arranques, separados \SI{1}{\second},
repitiendo el PRESTART completo. Para interpretar la distribuci\'on medida se
usa $\Delta\phi=2\pi f\Delta t$: \SI{400}{\micro\second} equivalen a
\SI{7.2}{\degree} a \SI{50}{\hertz} y \SI{14.4}{\degree} a
\SI{100}{\hertz}; un l\'imite de \SI{5}{\degree} exige
\SI{278}{\micro\second} y \SI{139}{\micro\second}, respectivamente.

\begin{figure}[!ht]
\centering
\includegraphics[width=0.86\textwidth]{sincronizacion_fase.png}
\caption{Presupuesto de sincronizaci\'on expresado en la variable que importa
para MASW: fase. La observaci\'on menor a \SI{400}{\micro\second} es preliminar
y no sustituye la distribuci\'on de \texttt{SCOPE\_MULTI}.}
\label{fig:sync-fase-hist}
\end{figure}

\textbf{Interfaz local en el nodo} (OLED por SPI $+$ botones). Cada esclavo
puede operarse sin maestro y sin tel\'efono. Junto con los comandos USB de la
Etapa~6, completa tres v\'ias independientes de operaci\'on:
interfaz local, USB y web.

\textbf{\texttt{PGAout} en toda la cadena de control}: comando \texttt{0xA8},
configuraci\'on dirigida por ESP-NOW, comando USB \texttt{pgaout N}, fila propia
en el panel web ---visible solo cuando el PSoC reporta
\texttt{hw\_class=GEO}--- con persistencia, y exportaci\'on como
\texttt{pgaout\_code}/\texttt{pgaout\_gain} en los metadatos.

\seleccionado{Que un par\'ametro nuevo aparezca simult\'aneamente en el firmware,
en los tres frentes de operaci\'on \emph{y} en los metadatos exportados es la
se\~nal de que la arquitectura de configuraci\'on madur\'o: agregar una etapa ya
no es un parche, es rellenar un patr\'on.}

\subsection{PCBs}

El repositorio \texttt{PCBs} contiene el proyecto en \textbf{KiCad} y en
\textbf{JitX} (m\'as un \texttt{JitX\_para\_port}), con \emph{commits} hasta el
9 de agosto.

\pendiente{No hay bit\'acora de esta etapa y la decisi\'on KiCad frente a JitX
no est\'a registrada en ning\'un lado. \textbf{¿Cu\'al es el proyecto vigente y por
qu\'e se explor\'o el otro?}}

\subsection{Balance de la etapa}

\begin{center}\small
\begin{tabularx}{\textwidth}{@{}l X@{}}
\toprule
\textbf{Se qued\'o} & Placa propia; IDAC como referencia de calibraci\'on (LSB \SI{3.75}{\milli\volt}); enlace asim\'etrico UART bajada / I\textsuperscript{2}C subida; etapa \texttt{PGAout}; interfaz local OLED $+$ botones; \emph{sync} \'unicamente por ESP-NOW con flanco generado localmente \\
\addlinespace[3pt]
\textbf{Se descart\'o} & VDAC como referencia; TX de la UART del PSoC; \emph{sync} por cable entre maestro y esclavo (\texttt{SYNC\_IN\_PIN} y su ISR); pin de LED en el TopDesign \\
\addlinespace[3pt]
\textbf{Qued\'o abierto} & Rehacer el presupuesto de \emph{jitter} sin cable de \emph{sync}; recalibrar el lazo PI con la nueva resoluci\'on del IDAC; validar la placa contra la matriz E1--E18 de la Etapa~8; decidir KiCad frente a JitX \\
\bottomrule
\end{tabularx}
\end{center}

\section{Inventario de lo construido}
\label{sec:inventario}

Qu\'e existe hoy y d\'onde est\'a. Sirve para reencontrar una pieza sin tener que
recorrer los diez repositorios.

\begin{figure}[ht]
\centering
\includegraphics[width=\textwidth]{tablero_evidencia.png}
\caption{Inventario visual de lo demostrado y de sus l\'imites. Complementa el
inventario de archivos con el nivel de evidencia disponible.}
\end{figure}

\subsection{Firmware}

\begin{center}\small
\begin{tabularx}{\textwidth}{@{}l X@{}}
\toprule
\multicolumn{2}{@{}l@{}}{\textbf{PSoC 5LP} --- \texttt{src/firmware/psoc}, \textasciitilde{}1200 archivos \texttt{.c}/\texttt{.h} (mayormente generados), 10 fuentes Verilog} \\
\midrule
Cadena anal\'ogica & PGA $\to$ BP $\to$ SUM $\to$ LP, con referencias por IDAC (antes VDAC) y AMux dedicados por etapa \\
Adquisici\'on & ADC $\Delta\Sigma$ multi-rango, \texttt{Filter} FIR de hardware, cadena DMA con canales crudo/filtrado $\times$ simple/lote \\
\texttt{superMaquina} & M\'aquina de estados de captura en UDB (Verilog): gobierna DRQ de DMA, inicio y fin de captura, y qu\'e IRQ llegan al ARM \\
Auto-calibraci\'on & Controlador PI con \emph{feedforward}, banda muerta por etapa, anti-\emph{windup} direccional, detector de \emph{lock}, persistencia en EEPROM con CRC-16 por c\'odigo de ganancia \\
Almacenamiento & FatFs sobre SD (\texttt{GEOLAST.BIN}, comandos \texttt{0xBC}--\texttt{0xBF}, banco \texttt{sdread N}) \\
Firmware unificado & Un \texttt{main.c} que compila como GEO o HAMMER seg\'un el TopDesign \\
\midrule
\multicolumn{2}{@{}l@{}}{\textbf{ESP32} --- \texttt{src/firmware/esp32}, maestro $+$ esclavos} \\
\midrule
Red & ESP-NOW para control y sincronizaci\'on; PRESTART/HOT\_WAIT; \emph{store-and-forward} \\
Enlace al PSoC & Asim\'etrico: UART en bajada, I\textsuperscript{2}C en subida (PSoC maestro) \\
Interfaz web & Aplicaci\'on servida por el maestro sobre su propio AP, con relay WebSocket \\
Interfaz local & OLED SPI $+$ botones en el esclavo (placa nueva) \\
Diagn\'ostico & Comandos USB \texttt{probe}, \texttt{status}, \texttt{stream}, \texttt{debugpsoc}, \texttt{pre}, \texttt{sync}, \texttt{startnow}, \texttt{cap}, \texttt{clear}, \texttt{stop}, \texttt{range N}, \texttt{pgaout N}, \texttt{sdread N} \\
Instrumentaci\'on & \texttt{CMD\_SCOPE\_START} y \texttt{SCOPE\_MULTI} para medir \emph{jitter} con osciloscopio \\
\bottomrule
\end{tabularx}
\end{center}

\subsection{Interfaces y post-proceso}

\begin{center}\small
\begin{tabularx}{\textwidth}{@{}l X@{}}
\toprule
\multicolumn{2}{@{}l@{}}{\textbf{Python} --- \texttt{src/interfaces/python} (68 m\'odulos) y \texttt{src/calculos\_modelados/python} (16)} \\
\midrule
Servidor & FastAPI con est\'aticos, pesta\~nas, tema, visor de se\~nal con decimado min/max \\
\texttt{review\_field\_data.py} & Revisi\'on de campo: orden por pico a pico, filtros, \'enfasis, alineaci\'on entre tandas, \emph{waterfall}, $f$--$k$, \emph{picking} de dispersi\'on, polaridad, dedup, mezcla de $F_s$ \\
\texttt{masw\_dispersion.py} & \emph{Gather} $\to$ imagen de dispersi\'on por \emph{phase-shift}; grupos de dispersi\'on multimodo con combinaci\'on ponderada \\
\texttt{masw\_inversion.py} & Fachada de inversi\'on a $V_S(z)$ con \emph{backends} seleccionables: \texttt{evodcinv}/\texttt{disba}, Monte Carlo, MASWavesPy, ADsurf y exportaci\'on a Geopsy \\
\texttt{deletion.py} & Cuarentena reversible de registros \\
\texttt{AnalysisJobs} & Cola persistente de an\'alisis \\
\texttt{smoke\_test.py} & \emph{Gate} de 58 verificaciones \\
C\'alculo de compensadores & Siete variantes de \texttt{calculoCompensador*}, \texttt{calculoMFB\_LPF.py}, \texttt{sim\_compensador\_bp.py} \\
\texttt{leva\_martillo/} & Modelo 2D del martillo de leva \\
\midrule
\multicolumn{2}{@{}l@{}}{\textbf{MATLAB} --- \texttt{src/interfaces/matlab} (14 \texttt{.m}) y \texttt{src/calculos\_modelados/matlab} (28)} \\
\midrule
\texttt{geophone\_scope.m} / \texttt{\_simple.m} & Osciloscopio virtual sobre puerto serie (657 l\'ineas la versi\'on completa) \\
\texttt{analisis\_espectral.m} & PSD, SNR por sesi\'on y espectrogramas por entidad (264 l\'ineas) \\
\texttt{AnalisisCircuito/} & \path{analizar_barridos_cadena_analogica.m}, \path{cancelar_armonicos_linea_ls.m}, \path{verificar_calibracion_potenciometro.m}, gr\'aficos GEO$\to$ADC \\
\texttt{SM-24/} & \texttt{Compensador.m}, \texttt{compensadorKai.m}, \texttt{compensadorPazos.m}, \texttt{filtroFir.m} \\
\texttt{CalculosRapidos/} & \texttt{ComparacionFiltros.m} (el gr\'afico del pivote), \texttt{notch.m}, \texttt{removeLineNoise\_SpectrumEstimation.m} \\
\texttt{pruebas\_teoricas/} & \texttt{p01}--\texttt{p05}: validaci\'on num\'erica de la teor\'ia \\
\texttt{Fisica loca/} & Martillo de leva en Simscape Multibody \\
\texttt{third-party/} & Librer\'ia MASW-Matlab, usada como referencia de contraste \\
\bottomrule
\end{tabularx}
\end{center}

\subsection{Datos}

\begin{center}\small
\begin{tabularx}{\textwidth}{@{}l r X@{}}
\toprule
\textbf{Conjunto} & \textbf{Tama\~no} & \textbf{Contenido} \\
\midrule
\texttt{raw/Canchiga} & 12{,}4~GiB & 74 capturas, 18 posiciones \\
\texttt{raw/Canchita} & 10{,}8~GiB & 87 capturas, 20 posiciones \\
\texttt{raw/Canchita\_2} & 5{,}0~GiB & 13 capturas, 10 posiciones \\
\texttt{raw/Estacionamiento} & 3~MiB & Prueba inicial en exteriores \\
\texttt{raw/Osciloscopio*} & \textasciitilde{}20~MiB & Barridos de la cadena anal\'ogica: tanda descartada del 20/07, calibraci\'on y verificaci\'on del 21/07 \\
\texttt{processed/Canchita\_procesado} & 530~MiB & \emph{Waterfall}, curva de dispersi\'on, perfil $V_S$, promedios, manifiesto \\
\texttt{processed/Canchita\_grupo1\_procesado} & 415~MiB & Grupo 1 procesado por separado \\
\bottomrule
\end{tabularx}
\end{center}

Todo versionado con Git LFS de almacenamiento por carpeta y deduplicaci\'on
SHA-256, con manifiesto de ruta, tama\~no y \emph{hash} por archivo.

\subsection{Documentaci\'on e investigaci\'on}

\begin{center}\small
\begin{tabularx}{\textwidth}{@{}l X@{}}
\toprule
Bit\'acora & 66 entradas diarias con an\'alisis de \emph{diff} \emph{commit} por \emph{commit}, \textasciitilde{}1{,}3~MB \\
Vault de investigaci\'on & 232 notas: 8 cap\'itulos de Foti resumidos $+$ \textasciitilde{}150 conceptos at\'omicos enlazados \\
Base bibliogr\'afica & 65 papers con DOI, categor\'ia y relevancia declarada \\
Art\'iculo URUCON & \emph{Automatic DC-Offset Calibration for a Low-Cost Geophone Front End}, 5 p\'aginas, con revisi\'on adversarial simulada \\
Protocolos & \texttt{PROTOCOLO\_BARRIDOS\_} \texttt{CADENA\_ANALOGICA.md}, \texttt{BUILD\_PROGRAM\_PSOC}, \texttt{PORT\_PLAN}, \texttt{PLAN\_CONECTIVIDAD\_} \texttt{MASTER.md} \\
\emph{Handoffs} & \texttt{HANDOFF\_SESION\_2026-07-12.md} (auditor\'ia pre-campo), \emph{handoff} de \emph{jitter} del 10/06 \\
Legado & \texttt{docs/legacy/pyqt\_field\_review/} --- \emph{snapshot} congelado de la app PyQt como referencia de paridad \\
\bottomrule
\end{tabularx}
\end{center}

\subsection{Instrumentaci\'on de laboratorio}

\begin{center}\small
\begin{tabularx}{\textwidth}{@{}l X@{}}
\toprule
Generador Agilent 33220A & Barridos logar\'itmicos con \texttt{SYNC} hacia el trigger externo \\
Osciloscopio Tektronix TDS2004B & 4 canales, uno por etapa anal\'ogica; sondas en $10\times$ obligatorio \\
Ge\'ofono SM-24 & \SI{10}{\hertz}, $\zeta_0 = 0{,}25$ \\
Martillo de impacto con sensor de fuerza PCB Piezotronics 086D20 & Canal de martillo \\
CY8CKIT-059 $+$ ESP32 DevKitC & \emph{Carrier} intermedia, previa a la placa propia \\
pyOCD & Depuraci\'on en placa (con \'el se vio la tormenta de IRQ) \\
\bottomrule
\end{tabularx}
\end{center}

\section{Tabla maestra de decisiones}
\label{sec:descartes}

Todo lo que se prob\'o y no qued\'o, con el motivo y qu\'e lo reemplaz\'o. Es la
secci\'on m\'as \'util para recordar \emph{por qu\'e} el sistema es como es ---y para
no volver a probar algo que ya se descart\'o con raz\'on.

\begin{figure}[ht]
\centering
\includegraphics[width=0.96\textwidth]{comparacion_metodos.png}
\caption{Contexto visual para las decisiones metodol\'ogicas: qu\'e entrega cada familia
y qu\'e costo operativo introduce.}
\end{figure}

\subsection{Cadena anal\'ogica}

\begin{center}\footnotesize
\begin{tabularx}{\textwidth}{@{}l X X l@{}}
\toprule
\textbf{Descartado} & \textbf{Motivo} & \textbf{Reemplazo} & \textbf{Fecha} \\
\midrule
Filtro RC pasivo de entrada
& Fase no lineal (descalificante para MASW); tolerancia de componentes; deriva t\'ermica; el PSoC ya tiene DSP
& Cadena FIR de fase lineal con decimaci\'on en cascada $\times4^3$
& 18/04 \\
\addlinespace[2pt]
Observador de estado (Simulink)
& Explorado y abandonado junto con el RC
& Compensaci\'on activa en anal\'ogico
& 18/04 \\
\addlinespace[2pt]
Topolog\'ia Sallen-Key / VCVS
& Realimentaci\'on positiva amplifica el error de componentes; riesgo de oscilaci\'on con $G>1$; impedancia de salida real. \textbf{No fue por respuesta}: ambas cumpl\'ian dentro de \SI{0.3}{\decibel}
& \textbf{MFB} (\emph{Multiple Feedback})
& 30/04 \\
\addlinespace[2pt]
Cancelaci\'on polo-cero como estrategia
& No es lo que hace falta: se busca \emph{elevar el amortiguamiento} manteniendo $\omega_n$
& Estrategia de Ma \emph{et al.}: pasa-todo menos pasa-banda
& abr--may \\
\addlinespace[2pt]
Amplificador operacional ideal en el modelo
& El dise\~no ``andaba en el papel'' pero no contemplaba $A_0$ ni GBW finitos
& Modelo no ideal con $R_{in}$, $R_{out}$ y $C_{in}$; la relaci\'on $f_p=f_t/A_0$ ya se corrigi\'o, pero la s\'intesis y la medici\'on deben repetirse antes de reutilizar el ajuste
& 17/05 \\
\addlinespace[2pt]
Discretizar a valores comerciales antes de optimizar
& Deja el \'optimo fuera de la grilla
& Optimizar en el continuo $\to$ restringir por f\'isica $\to$ \emph{snap} y redes en T al final
& 18/05 \\
\addlinespace[2pt]
Rama BP con \SI{8200}{\ohm} (dise\~no pide \SI{7501.9}{\ohm})
& El compensador realizaba $\zeta = 83{,}661$ en vez de $0{,}25$
& \SI{6.8}{\kilo\ohm} $+$ potenci\'ometro de \SI{2}{\kilo\ohm}
& 21/07 \\
\addlinespace[2pt]
Sondas de osciloscopio en $1\times$
& Su capacitancia es comparable al $C_f = \SI{177}{\pico\farad}$ del nodo BP: el instrumento carga el circuito que mide
& Sondas compensadas en $10\times$
& 20/07 \\
\addlinespace[2pt]
VDAC como referencia de calibraci\'on
& Cambio de placa
& IDAC, LSB \SI{3.75}{\milli\volt}
& 03/08 \\
\bottomrule
\end{tabularx}
\end{center}

\subsection{Procesamiento de se\~nal}

\begin{center}\footnotesize
\begin{tabularx}{\textwidth}{@{}l X X l@{}}
\toprule
\textbf{Descartado} & \textbf{Motivo} & \textbf{Reemplazo} & \textbf{Fecha} \\
\midrule
Filtro adaptativo LMS para \SI{50}{\hertz}
& Tiene transitorio de convergencia, que en una captura de pocos segundos ocupa fracci\'on apreciable del registro; y su comportamiento depende de la se\~nal, lo que rompe la comparabilidad entre canales
& Estimaci\'on directa por m\'inimos cuadrados
& 07/06 \\
\addlinespace[2pt]
Notch de \SI{50}{\hertz} (FIR y notch arm\'onico en el cliente)
& La \SI{50}{\hertz} cae \emph{dentro} de la banda \'util del MASW: el notch borra se\~nal e introduce retardo de grupo justo donde se mide fase
& B\'usqueda acotada de frecuencia $+$ regresi\'on arm\'onica lineal en post-proceso: sin retardo y con el crudo preservado; puede atenuar contenido coherente con las bases ajustadas
& 09/06 \\
\addlinespace[2pt]
$F_s$ \'unica fija
& Distintos compromisos entre resoluci\'on temporal y volumen por captura
& En las campa\~nas hist\'oricas se usaron $2929$ y \SI{1020}{\hertz}, con mezcla posterior por \texttt{resample\_poly}. El firmware actual usa \SI{2604}{\hertz} nativos a 18 bits y entrega $F_s=2604/N$
& 07/07 \\
\addlinespace[2pt]
Rango de ADC \'unico
& El nivel de se\~nal var\'ia mucho con la distancia fuente--receptor
& 2 y luego 4 rangos seleccionables desde la web, validados con \texttt{range N}
& 07--08/07 \\
\bottomrule
\end{tabularx}
\end{center}

\subsection{Firmware y arquitectura embebida}

\begin{center}\footnotesize
\begin{tabularx}{\textwidth}{@{}l X X l@{}}
\toprule
\textbf{Descartado} & \textbf{Motivo} & \textbf{Reemplazo} & \textbf{Fecha} \\
\midrule
Firmware bifurcado GEO / HAMMER
& Dos implementaciones del mismo protocolo divergen y desincronizan el arreglo
& \'Unico \texttt{main.c} con detecci\'on de hardware en tiempo de compilaci\'on
& 04/06 \\
\addlinespace[2pt]
Bisecci\'on binaria sobre c\'odigos VDAC
& Ciega: no aprovecha que la planta es conocida; una versi\'on por etapa
& Controlador PI unificado
& 23/06 \\
\addlinespace[2pt]
Servo lento
& Convergencia por integraci\'on pura
& Controlador PI unificado
& 23/06 \\
\addlinespace[2pt]
T\'ermino derivativo del PID
& Amplifica el ruido de medici\'on del ADC sin aportar en planta cuasi-est\'atica (\emph{``overfitting al ruido''})
& \textbf{PI} con \emph{feedforward} f\'isico
& 23/06 \\
\addlinespace[2pt]
Anti-\emph{windup} sim\'etrico
& Congelar el integrador en ambos sentidos deja al lazo sin poder salir del riel
& Anti-\emph{windup} \textbf{direccional}
& 25/06 \\
\addlinespace[2pt]
ISR cl\'asica de adquisici\'on
& El ARM en el camino de la muestra introduce \emph{jitter}
& Cadena DMA con \texttt{Filter} FIR de hardware
& 17/06 \\
\addlinespace[2pt]
\texttt{Reg\_Select} como gobierno de la captura
& El temporizado heredaba el \emph{jitter} de las interrupciones. \emph{Determinismo temporal es exactitud de medici\'on}
& \texttt{superMaquina}: m\'aquina de estados en UDB (Verilog)
& 02/07 \\
\addlinespace[2pt]
\emph{Tick} de sistema de \SI{10}{\milli\second}
& Un \'unico \emph{tick} para funciones de per\'iodos muy distintos
& Tres \emph{timers} fijos dedicados
& 02/07 \\
\addlinespace[2pt]
Interrupci\'on por \texttt{TC}
& Pol\'itica de determinismo: la ISR maneja \emph{solo} muestras
& \emph{Polling} de \texttt{STATUS\_TC}; \texttt{CLEAR\_FLAGS} prohibido en \texttt{SAMPLING}
& 07/07 \\
\bottomrule
\end{tabularx}
\end{center}

\subsection{Comunicaci\'on y sincronizaci\'on}

\begin{center}\footnotesize
\begin{tabularx}{\textwidth}{@{}l X X l@{}}
\toprule
\textbf{Descartado} & \textbf{Motivo} & \textbf{Reemplazo} & \textbf{Fecha} \\
\midrule
Espera fija tras \texttt{ARM}
& El esclavo pod\'ia recibir \texttt{START} sin estar en \texttt{ARMED}: el nodo simplemente no arranca, y el resto del arreglo s\'i
& \emph{Probe}/ACK con RTT medido, y luego PRESTART con HOT\_WAIT
& 25--26/05 \\
\addlinespace[2pt]
Pulsos de \emph{debug} con \texttt{delayMicroseconds()}
& Bloquean \SI{1}{\milli\second} y hacen perder callbacks de ESP-NOW: \emph{el instrumento de medici\'on causaba la falla que se buscaba}
& Pulsos no bloqueantes por marca de tiempo
& 26/05 \\
\addlinespace[2pt]
Enlace SPI PSoC$\leftrightarrow$ESP
& Reemplazado en el refactor del PR \#2
& UART, y despu\'es I\textsuperscript{2}C en la subida
& 31/05 \\
\addlinespace[2pt]
ACK emitido antes del flanco f\'isico
& El maestro cre\'ia sincronizado un arreglo que a\'un no lo estaba; se manifestaba como \emph{jitter}, no como falla
& Flanco local inmediato y silencio de radio durante el muestreo; ACK s\'olo para sondeo, rechazo o finalizaci\'on
& 10/06 \\
\addlinespace[2pt]
UART de \SI{115200}{} como subida de muestras
& Cuello de botella del \emph{stream}
& I\textsuperscript{2}C con PSoC de maestro; \emph{mismo protocolo, otro medio}
& 03/08 \\
\addlinespace[2pt]
\emph{Sync} por cable maestro$\to$esclavo (\texttt{SYNC\_IN\_PIN})
& Un cable menos por nodo en campo
& Disparo \'unicamente por ESP-NOW; el esclavo genera el flanco local
& 03/08 \\
\bottomrule
\end{tabularx}
\end{center}

\subsection{Robustez (auditor\'ia pre-campo)}

\begin{center}\footnotesize
\begin{tabularx}{\textwidth}{@{}l X X l@{}}
\toprule
\textbf{Descartado} & \textbf{Motivo} & \textbf{Reemplazo} & \textbf{Fecha} \\
\midrule
\emph{Watchdog} \texttt{SYNC\_STALL} con margen fijo
& \textbf{F1:} en capturas largas disparaba solo y \emph{re-disparaba} \texttt{START\_NOW}, destruyendo la sesi\'on. 10 minutos $=$ 0 muestras
& Margen proporcional a la duraci\'on y aborto limpio
& 11/07 \\
\addlinespace[2pt]
Reintento de PRESTART sin cubrir el ARMED silencioso del PSoC
& \textbf{F9:} START romp\'ia el nodo hasta reset f\'isico, siempre
& Carrera cerrada
& 12/07 \\
\addlinespace[2pt]
Herramienta WS sin \emph{watchdog} de silencio
& \textbf{F5:} colgada indefinida ante socket medio abierto
& \emph{Watchdog} de silencio $+$ endurecimiento pre-DATA
& 12/07 \\
\addlinespace[2pt]
Criterios de prueba bloqueados por contadores benignos
& \textbf{F8:} el contador \texttt{drop} del \emph{framer} bloqueaba pruebas sin ser un error
& Solo \texttt{bBad}/\texttt{badLen} bloquean
& 12/07 \\
\bottomrule
\end{tabularx}
\end{center}

\subsection{Herramientas e infraestructura}

\begin{center}\footnotesize
\begin{tabularx}{\textwidth}{@{}l X X l@{}}
\toprule
\textbf{Descartado} & \textbf{Motivo} & \textbf{Reemplazo} & \textbf{Fecha} \\
\midrule
Repositorio monol\'itico
& Ciclos de vida distintos; \textasciitilde{}29~GiB de datos crudos envenenaban el historial
& Superproyecto con subm\'odulos $+$ LFS con dedup SHA-256
& 21/07 \\
\addlinespace[2pt]
\texttt{src/} organizado por lenguaje
& Separaba cosas que se usan juntas: hay interfaces y modelado en dos lenguajes cada uno
& Organizaci\'on \textbf{por prop\'osito}
& 24/07 \\
\addlinespace[2pt]
\texttt{http.server} artesanal
& L\'imites para servir la aplicaci\'on de revisi\'on
& FastAPI $+$ est\'aticos
& 25/07 \\
\addlinespace[2pt]
App de escritorio PyQt \texttt{field\_review\_app}
& Se port\'o a web. \textbf{No se borr\'o}: se congel\'o como referencia de paridad en \texttt{docs/legacy/}
& Interfaz web con paridad funcional
& 26--29/07 \\
\addlinespace[2pt]
Borrado f\'isico de registros rechazados
& Un descarte err\'oneo de datos de campo es irreparable
& Cuarentena reversible (\texttt{deletion.py})
& 29/07 \\
\addlinespace[2pt]
\emph{Sketch} \texttt{MasterFoo}
& Descartable por dise\~no: exist\'ia para responder si serv\'ia el \emph{hotspot} del celular. Respondi\'o
& Plan de conectividad modo ENLACE
& 24/07 \\
\addlinespace[2pt]
Tanda \texttt{Osciloscopio\_descartado\_2026-07-20\_BP\_43k}
& Rama BP con el resistor equivocado. \textbf{Se conserva con el motivo en el nombre}
& Tanda calibrada del 21/07
& 20/07 \\
\bottomrule
\end{tabularx}
\end{center}

\subsection{Los principios que sobrevivieron a todo}

\seleccionado{Ocho decisiones nunca se revirtieron, y son las que en la
pr\'actica definen el sistema:
\begin{enumerate}[itemsep=0.15em,topsep=0.25em]
\item \textbf{Fase lineal siempre.} Todo el filtrado que toca la se\~nal
conserva la fase, porque el m\'etodo mide fase.
\item \textbf{El instante cero lo define un flanco f\'isico}, no un mensaje.
\item \textbf{Determinismo temporal por hardware}, no por software.
\item \textbf{El protocolo es independiente del medio.} SPI $\to$ UART $\to$
I\textsuperscript{2}C sin tocar los \emph{frames}.
\item \textbf{El dato crudo no se toca.} Toda correcci\'on discutible ---red
el\'ectrica, polaridad, \'enfasis--- vive en post-proceso y es reversible.
\item \textbf{Nada se borra: se cuarentena o se archiva con el motivo.}
\item \textbf{Validar antes de borrar.} Lo reemplazado queda como referencia
hasta que el reemplazo se prueba.
\item \textbf{Un solo firmware, un solo controlador, un solo protocolo.}
Cada bifurcaci\'on evitada es un modo de falla evitado.
\end{enumerate}}

\section{Qu\'e queda abierto}
\label{sec:pendientes}

Ordenado por impacto sobre el objetivo del proyecto, no por dificultad.

\begin{figure}[ht]
\centering
\includegraphics[width=\textwidth]{roadmap_validacion.png}
\caption{Hoja de ruta cualitativa por dependencia y valor de evidencia. No reemplaza un
cronograma acordado con tutor\'ia.}
\end{figure}

\subsection{Lo que hay que resolver antes de escribir la entrega}

\begin{enumerate}[itemsep=0.45em]

\item \textbf{Recalcular y medir el modelo no ideal del operacional.} Los
scripts usaban $A_0=\SI{90}{\decibel}$ y \SI{8}{\mega\hertz} como polo dominante;
si \SI{8}{\mega\hertz} es GBW, el polo monopolo es aproximadamente
\SI{253}{\hertz}. La funci\'on de lazo abierto ya se corrigi\'o mediante
$f_p=f_t/A_0$. Falta identificar el operacional configurado y repetir la
s\'intesis y el contraste experimental.

\item \textbf{La profundidad alcanzable no est\'a sostenida por la geometr\'ia}
(Etapa~7). Con apertura $L = \SI{40}{\meter}$, el soporte geom\'etrico
llega a \textasciitilde{}\SI{20}{\meter}, pero el modelo invertido se extiende
hasta \SI{82.6}{\meter}. O se alarga el arreglo, o se declara expl\'icitamente
que por debajo de \SI{20}{\meter} el perfil es extrapolaci\'on del modelo. Para
el objetivo de \SI{50}{\meter} no alcanza con la banda: hace falta apertura.

\item \textbf{La mitad lejana del arreglo est\'a limitada por ruido}
(Etapa~7): SNR mediana \SI{4.9}{\decibel} entre \SI{30}{} y
\SI{50}{\meter}, con posiciones bajo \SI{0}{\decibel}. Alargar el arreglo sin
aumentar la energ\'ia de la fuente no resuelve nada. Esto le da prioridad al
martillo de leva.

\item \textbf{La muesca de compensaci\'on est\'a \SI{26.33}{\decibel} por encima
de su profundidad de dise\~no} (Etapa~9), y la causa es que el ajuste
requerido es de \SI{1.98}{\ohm} sobre un potenci\'ometro de \SI{2}{\kilo\ohm}
---una parte en mil. La acci\'on registrada es recalibrar alrededor de
\SI{10.2046}{\hertz} con mayor sensibilidad vertical en CH3 y repetir un
barrido estrecho; la acci\'on de fondo es rehacer el divisor con resoluci\'on
adecuada.

\end{enumerate}

\subsection{Mediciones que faltan}

\begin{center}\small
\begin{tabularx}{\textwidth}{@{}X l@{}}
\toprule
\textbf{Medici\'on} & \textbf{Impacto} \\
\midrule
\emph{Jitter} real del arreglo completo desplegado. Adem\'as, tras eliminar el \emph{sync} por cable (Etapa~11), \textbf{el presupuesto de \emph{jitter} hay que rehacerlo entero}: ahora todo el error de sincronizaci\'on pasa por el enlace de radio & Alto \\
\addlinespace[2pt]
Apareamiento entre canales. Un desfasaje entre canales se traduce directamente en error de velocidad de fase & Alto \\
\addlinespace[2pt]
Banda plana efectivamente lograda tras compensar, medida. Es el resultado que cierra la l\'inea del \emph{front-end} y permite compararse con los \SIrange{1}{100}{\hertz} de Ma \emph{et al.} & Alto \\
\addlinespace[2pt]
Ruido del sistema referido a la entrada, contra el l\'imite t\'ermico del ge\'ofono y separado por etapa. La compensaci\'on modifica la respuesta al ruido y debe caracterizarse de extremo a extremo & Medio-alto \\
\addlinespace[2pt]
Par\'ametros medidos del SM-24 usado: frecuencia natural, amortiguamiento y sensibilidad. La hoja de datos especifica resistencia de bobina, masa m\'ovil y sensibilidad nominal, pero no informa inductancia & Medio \\
\addlinespace[2pt]
Autonom\'ia del nodo y consumo. No registrado en ning\'un lado & Medio \\
\addlinespace[2pt]
Validaci\'on externa del perfil $V_S$ contra SPT o ensayo invasivo. Sin ella, el resultado demuestra que la cadena cierra, no que el perfil sea correcto & Alto \\
\addlinespace[2pt]
Recalibrar el lazo PI con la resoluci\'on del IDAC (LSB \SI{3.75}{\milli\volt}) y revalidar la placa nueva contra la matriz E1--E18 & Alto \\
\bottomrule
\end{tabularx}
\end{center}

\subsection{Preguntas sin respuesta registrada}

\pendiente{\textbf{KiCad o JitX.} El repositorio \texttt{PCBs} tiene los dos
proyectos y no hay bit\'acora de agosto. \textbf{¿Cu\'al es el vigente y por qu\'e
se explor\'o el otro?}}

\pendiente{\textbf{Adquisici\'on pasiva.} El sistema puede, en principio,
operar en modo pasivo (ReMi, SPAC) simplemente capturando sin disparar la
fuente ---la captura continua de 10 minutos ya est\'a validada. Combinar activo
y pasivo extiende la curva de dispersi\'on hacia frecuencias bajas, es decir
hacia mayores profundidades, sin hardware adicional. Dado que la apertura es
hoy la restricci\'on activa, es la extensi\'on de alcance m\'as barata disponible.
\textbf{¿Se consider\'o?}}

\pendiente{\textbf{Los nombres internos de los sitios} (Canchiga, Canchita,
Estacionamiento) ¿se mantienen o se formalizan en la entrega?}

\pendiente{\textbf{El uso de asistencia de IA} fue intensivo y est\'a
documentado en detalle en la bit\'acora. ¿Se explicita en la entrega, y con qu\'e
encuadre metodol\'ogico? Conviene decidirlo antes de que lo pregunte la mesa.}

\subsection{Deuda t\'ecnica conocida}

\begin{center}\small
\begin{tabularx}{\textwidth}{@{}l X@{}}
\toprule
Vault de investigaci\'on & No se audita desde abril. Puede tener conceptos escritos suponiendo compensaci\'on por cancelaci\'on polo-cero, enfoque abandonado \\
\addlinespace[2pt]
H4 & \emph{Prefetch} de bloque en el PSoC dar\'ia un volcado \textasciitilde{}$2\times$ m\'as r\'apido. Identificado, no implementado \\
\addlinespace[2pt]
W3 & STOP durante el volcado deja los datos en la SD sin forma de re-volcarlos desde la interfaz \\
\addlinespace[2pt]
F2 / W5 & Cosm\'eticos, documentados \\
\addlinespace[2pt]
Nombres de carpeta de campo & No siempre coinciden con la posici\'on declarada en el metadato (\texttt{muestra28\_22} declara \SI{30}{\meter}). \textbf{El metadato es el autoritativo}; conviene no volver a leer geometr\'ia del nombre \\
\bottomrule
\end{tabularx}
\end{center}

\subsection{Nota de cierre}

Vale la pena registrar el balance sin cerrar antes de tiempo lo que a\'un exige
medici\'on. El proyecto demostr\'o la cadena completa, una compensaci\'on
anal\'ogica repetible aunque todav\'ia lejos de la cancelaci\'on fina de dise\~no,
auto-calibraci\'on con evidencia estad\'istica y sincronizaci\'on inal\'ambrica
funcional en banco. En campo aparecieron adem\'as dos l\'imites independientes:
la apertura del arreglo y la energ\'ia de la fuente. La contribuci\'on est\'a
sostenida; la banda efectiva, el \emph{jitter} extremo a extremo y la validez a
\SI{50}{\meter} siguen siendo resultados por cerrar, no logros demostrados.


\end{document}
